\chapter{Mixture of Experts：参数效率的极致}
\label{ch:moe}
%% ============================================================

\section{核心思想：不是每个神经元都需要工作}

想象一个公司有 256 名专家，分别精通不同领域。
当一个任务到来时，你不需要让所有 256 人都参与，
你只需要挑选最相关的 8 个人来处理。
这就是 Mixture of Experts (MoE) 的核心思想。

\subsection{历史脉络}

MoE 的思想远比深度学习更古老。
1991 年，Jacobs 等人提出了最早的 MoE 框架，% NEW REF: jacobs1991moe = Adaptive Mixtures of Local Experts, Jacobs et al., Neural Computation 1991
用多个「专家」网络分别处理输入空间的不同区域，
一个「门控」网络决定每个输入由哪些专家处理。
但这个想法在之后的 25 年里基本停留在学术论文中，
没有被大规模应用。

2017 年，Shazeer 等人~\cite{shazeer2017moe}
将 MoE 引入了 NLP 领域，
在 LSTM 语言模型中嵌入了稀疏门控的 MoE 层，
用 $>$1000 个专家实现了当时的最大模型（137B 参数）。
这篇论文的标题直接叫做
“Outrageously Large Neural Networks”，
「大得离谱的神经网络」。

2022 年，Switch Transformer~\cite{fedus2022switch}
将 MoE 扩展到 Transformer 架构，
用 top-1 路由（每个 token 只发给 1 个专家）
训练了一个 1.6T 参数的模型。
Switch Transformer 证明了一个关键点：
\textbf{即使只激活一个专家，
只要专家数量足够多、路由足够准确，
模型性能就可以超越同规模的稠密模型}。

自此，MoE 从一种学术好奇心
变成了 LLM 训练的主流架构选择。

值得注意的是，MoE 思想的发展并非一帆风顺。
在 Shazeer 2017 年的工作与 Switch Transformer 2022 年之间，
有一段“MoE 低谷期”，
多个研究团队尝试将 MoE 应用到 Transformer 上但遇到了严重的训练不稳定问题。
GShard~\cite{lepikhin2021gshard}（2021）是这段时期的一个重要突破，
它在机器翻译任务上成功训练了 600B 参数的 MoE Transformer，
证明了大规模 MoE Transformer 的可行性。
GShard 引入了两个至今仍被广泛使用的关键技术：
\textbf{容量因子}（capacity factor）限制每个专家处理的最大 token 数，
以及\textbf{随机路由}作为 Top-2 路由的第二选择
（第一个专家按概率选择，第二个专家按残余概率随机选择）。
这些看似简单的工程技巧解决了当时 MoE 训练中
最棘手的负载不均衡和训练崩溃问题。

\subsection{稠密模型 vs 稀疏模型}

传统的「稠密」模型（Dense Model）在处理每个 token 时
使用所有参数，参数量直接决定了计算量。
MoE 模型则打破了这种绑定：
\textbf{总参数量决定了模型的知识容量，
但每个 token 的计算量只取决于被激活的参数}。

DeepSeek-V3 的数据最能说明这种效率：
总参数 671B（知识容量巨大），
每个 token 仅激活 37B（计算成本低），
激活率约 5.5\%。

用一个类比来理解：
稠密模型像一个通才，不管什么问题，
他调动全部知识和精力来回答，效果不错但很累。
MoE 模型像一个高效的团队，
不管什么问题，先判断这是什么类型的问题，
然后把最擅长的几个人召集起来处理，
其他人可以休息。
团队的总知识量（总参数）可以很大，
但每次实际消耗的精力（计算量）很小。

让我们用具体的数字来感受这种效率。
下表展示了几个代表性模型的参数量和推理 FLOP 对比：

\begin{center}\small
\renewcommand{\arraystretch}{1.3}
\begin{tabular}{lrrrr}
\toprule
\rowcolor{figLightGray}
\textbf{模型} & \textbf{总参数} & \textbf{活跃参数} & \textbf{每 token FLOP}
& \textbf{效率比} \\
\midrule
LLaMA 3 70B（Dense） & 70B & 70B & $\sim$140 GFLOP & 1.0$\times$ \\
Mixtral 8x7B (MoE)   & 47B & 13B & $\sim$26 GFLOP  & 5.4$\times$ \\
DeepSeek-V3 (MoE)    & 671B & 37B & $\sim$74 GFLOP & 1.9$\times$ \\
Qwen3-Next (MoE)     & 80B & 3B  & $\sim$6 GFLOP   & 23.3$\times$ \\
\bottomrule
\end{tabular}
\end{center}

效率比以 LLaMA 3 70B 为基线。
Mixtral 以不到 70B 的总参数实现了 5.4 倍的推理效率提升，
同时性能匹配或超越 LLaMA 2 70B。
DeepSeek-V3 虽然总参数是 70B 的 9.6 倍，
但推理 FLOP 只有约一半，
用更多的“总知识”换取了更低的“单次计算成本”。
Qwen3-Next 将这一逻辑推到了极端：
80B 的模型只用 3B 的计算量，效率比达到 23.3 倍。

\section{MoE 架构总览}

在深入路由机制之前，
先从整体上理解 MoE 层在 Transformer 中的位置和数据流。
图~\ref{fig:moe-architecture} 展示了一个 MoE 层的完整结构。

\begin{figure}[htbp]
\centering
\begin{tikzpicture}[
  >=Stealth,
  node distance=8mm and 12mm,
  every node/.style={font=\small},
  box/.style={draw, rounded corners=3pt, minimum height=10mm,
              minimum width=25mm, inner sep=8pt, align=center},
  expert/.style={box, fill=figBlue!15, draw=figBlue!60},
  router/.style={box, fill=figOrange!20, draw=figOrange!70},
  shared/.style={box, fill=figGreen!15, draw=figGreen!60},
  combine/.style={box, fill=figRed!12, draw=figRed!60},
  input/.style={box, fill=figLightGray, draw=figGray},
  arrow/.style={->, thick, figGray},
  dasharrow/.style={->, thick, dashed, figGray!60},
]

% Input
\node[input] (input) {Token $\mathbf{x}$\\(隐状态 $d$-维)};

% Router
\node[router, above=12mm of input] (router) {Router\\$W_r \mathbf{x} \to \mathbb{R}^N$};

% Top-K selection
\node[box, fill=figOrange!10, draw=figOrange!50,
      above=10mm of router, minimum width=35mm] (topk)
      {Top-K 选择 + Softmax\\$g_1, g_2, \ldots, g_K$};

% Experts - spread out horizontally
\node[expert, above left=14mm and 30mm of topk] (e1) {Expert 1};
\node[expert, above left=14mm and 0mm of topk] (e2) {Expert $i$};
\node[expert, above right=14mm and 30mm of topk] (eN) {Expert $N$};

% Dots between experts
\node[above right=14mm and 12mm of topk, font=\large] {$\cdots$};

% Shared expert (on the right side)
\node[shared, right=25mm of topk] (shared) {Shared\\Expert};

% Combine
\node[combine, above=55mm of topk] (combine)
     {加权求和\\$\mathbf{y} = \sum g_i \cdot E_i(\mathbf{x}) + E_s(\mathbf{x})$};

% Output
\node[input, above=10mm of combine] (output)
     {Output $\mathbf{y}$\\(残差连接后)};

% Arrows: input -> router -> topk
\draw[arrow] (input) -- (router);
\draw[arrow] (router) -- (topk);

% Arrows: topk -> experts (selected)
\draw[arrow, figBlue] (topk.north) -- ++(0, 4mm) -| (e1.south)
     node[pos=0.75, left, font=\scriptsize, figBlue] {$g_1$};
\draw[arrow, figBlue] (topk.north) -- (e2.south)
     node[pos=0.5, right, font=\scriptsize, figBlue] {$g_i$};
\draw[arrow, figBlue] (topk.north) -- ++(0, 4mm) -| (eN.south)
     node[pos=0.75, right, font=\scriptsize, figBlue] {$g_K$};

% Arrow: input -> shared expert (bypass router, wider clearance for topk)
\draw[arrow, figGreen] (input.east) -- ++(14mm, 0) |- (shared.south)
     node[pos=0.25, right, font=\scriptsize, figGreen] {全部 token};

% Arrows: experts -> combine
\draw[arrow, figBlue] (e1.north) |- (combine.west);
\draw[arrow, figBlue] (e2.north) -- (combine.south);
\draw[arrow, figBlue] (eN.north) |- ([yshift=2mm]combine.east);

% Arrow: shared -> combine (offset to avoid overlap with eN arrow)
\draw[arrow, figGreen] (shared.north) |- ([yshift=-2mm]combine.east);

% Arrow: combine -> output
\draw[arrow] (combine) -- (output);

% Annotations
\node[right=3mm of router, font=\scriptsize, text=figGray, text width=28mm]
     {线性投影\\无偏置、无激活函数\\参数量 $\sim d \times N$};

% Inactive experts annotation
\draw[dasharrow] (topk.west) -- ++(-15mm, 0)
     node[left, font=\scriptsize, text=figGray!80, text width=22mm, align=right]
     {未选中的 $N{-}K$\\个专家不参与\\计算（零 FLOP）};

\end{tikzpicture}
\caption{MoE 层的架构示意图。输入 token 经过轻量级路由器选出 Top-K 个专家，
各专家的输出按路由权重加权求和。共享专家（DeepSeek 设计）处理所有 token 的通用知识，
不经过路由选择。未被选中的专家完全不参与计算，这是 MoE 稀疏激活的核心。}
\label{fig:moe-architecture}
\end{figure}

MoE 层替换了 Transformer 中的标准 FFN（前馈网络）层，
在多头注意力之后、残差连接之前。
每个 Transformer 块中可能有一个或多个 MoE 层，
不同模型的选择不同。
Mixtral 将每一层的 FFN 都替换为 MoE 层，
而一些模型（如早期的 GShard~\cite{lepikhin2021gshard}）
只在偶数层使用 MoE，奇数层保留标准 FFN，
后者的设计考量是减少 All-to-All 通信次数
和维持部分计算的确定性。

\section{路由机制：MoE 的核心决策}

对于输入 token $\mathbf{x}$，
路由器（一个简单的线性层）计算每个专家的分数，
然后选择得分最高的 $K$ 个专家：
\begin{equation}
  g_i = \frac{e^{s_i}}{\sum_{j \in \mathrm{TopK}} e^{s_j}}, \quad
  s = W_{\mathrm{router}} \mathbf{x}, \quad
  \mathbf{y} = \sum_{i \in \mathrm{TopK}} g_i \cdot \mathrm{Expert}_i(\mathbf{x})
  \label{eq:moe-routing}
\end{equation}

路由器本身极其轻量，
它只是一个 $[d, N_{\mathrm{experts}}]$ 的线性投影，
没有偏置，没有激活函数。
它的参数量微不足道（$4096 \times 256 \approx 1M$），
但它做出的决策至关重要：
将每个 token 送到最合适的专家去处理。

训练过程中，路由器通过梯度下降与专家一起优化。
一个有趣的现象是路由器会自发学到有意义的分工：
某些专家可能专注于代码相关的 token，
某些专注于数学符号，
某些专注于特定语言的文本。
这种\textbf{自动专业化}不是人为设定的：
它完全是训练过程中涌现出来的。

但这里有一个需要仔细理解的技术细节：
公式~\ref{eq:moe-routing} 中的 softmax 是在
\textbf{已选中的 Top-K 个专家}上计算的，
而不是在所有 $N$ 个专家上计算。
这意味着未被选中的专家不仅不参与前向计算，
它们的路由概率也不进入 softmax 归一化。
这种设计被称为 “softmax-then-top-k” 或 “top-k-then-softmax”，
取决于具体实现~\cite{shazeer2017moe}。
两种顺序的差异微妙但重要：
先 softmax 再选 top-k 会在梯度中引入所有专家的信号
（因为 softmax 的分母包含所有专家），
而先选 top-k 再 softmax 则只有被选中的专家收到梯度更新。
后者更高效但可能导致未被选中的专家永远得不到学习信号，
这是专家坍缩问题的一个来源。

从工程角度看，路由器虽然简单，
但它在 MoE 训练中的地位极其关键。
一个训练不好的路由器可以导致整个模型的失败，
即使你有 256 个强大的专家，
如果路由器无法正确匹配 token 和专家，
模型的表现可能还不如一个只有 10B 参数的稠密模型。
这就像一个优秀的调度系统对于出租车平台的重要性：
即使有 1000 辆车，如果调度不合理，
乘客等待时间可能比只有 100 辆车但调度优秀的平台更长。

\subsection{Top-K 路由：主流方案的深入分析}

Top-K 路由是目前最主流的方案。
$K$ 的选择涉及质量与效率的权衡：

\textbf{Top-1}（Switch Transformer~\cite{fedus2022switch}）：
每个 token 只发给 1 个专家。
计算效率最高，但单个专家可能无法充分处理复杂 token。
而且 top-1 的路由决策更「尖锐」，
一旦选错专家，没有其他专家来补救。

\textbf{Top-2}（Mixtral~\cite{jiang2024mixtral}、GShard）：
每个 token 发给 2 个专家，% NEW REF: lepikhin2021gshard = GShard: Scaling Giant Models with Conditional Computation, Lepikhin et al., ICLR 2021
结果按路由权重加权求和。
这是一种「保险」策略：
即使一个专家不太合适，
另一个通常能弥补。
Top-2 是最常见的选择，
在质量和效率之间取得了良好平衡。

\textbf{Top-8}（DeepSeek-V3~\cite{deepseekv3}）：
每个 token 发给 8 个专家。
看起来计算量更大？
不，因为 DeepSeek 使用了 256 个\textbf{细粒度小专家}，
每个专家只有标准专家的 $\sim$1/32 大小。
选 8 个小专家的总计算量与选 2 个大专家相当，
但组合空间爆炸式增长（从 28 种增加到 44 亿种）。

\subsection{Token Choice vs Expert Choice}

上述路由方式都是\textbf{Token Choice}，
每个 token 独立选择自己要去的专家。
这会导致一个固有问题：
热门专家可能收到过多 token，冷门专家闲置。

Zhou 等人（2022）提出了\textbf{Expert Choice} 路由：% NEW REF: zhou2022expertchoice = Mixture-of-Experts with Expert Choice Routing, Zhou et al., NeurIPS 2022
反过来，让专家选择 token。
每个专家从所有 token 中选取得分最高的 $k$ 个来处理。
这天然保证了负载均衡，
每个专家处理相同数量的 token，不存在过载或空闲。

但 Expert Choice 有一个问题：
某些 token 可能被多个专家选中（计算冗余），
而某些 token 可能不被任何专家选中（信息丢失）。
后者尤其危险，被“遗漏”的 token 在该 MoE 层的输出为零
（或仅来自残差连接），
相当于这一层对这些 token 没有做任何有意义的计算。
如果一个 token 在多个连续的 MoE 层中都被遗漏，
它的表示几乎得不到更新，
信息传播被严重阻断。

Expert Choice 的另一个限制是它对自回归推理不友好。
在训练中，Expert Choice 可以在整个 batch 上做全局的专家分配，
因为所有 token 都是可见的。
但在自回归推理的 decode 阶段，
每次只生成一个 token，
“让专家从所有 token 中选择”这个操作退化为
“让专家决定是否处理这一个 token”：
失去了全局优化的空间。
而且在 decode 阶段如果某个 token 没有被任何专家选中，
模型输出直接出错，无法通过后续 token 来弥补。
这就是为什么 Expert Choice 主要被研究社区关注，
而大规模生产模型几乎都使用 Token Choice。

在实践中，Token Choice + 负载均衡机制
仍然是更主流的方案。
但 Expert Choice 的思想并非没有价值，
它启发了一些混合策略：
在训练时用 Expert Choice 的负载均衡特性
作为辅助信号来优化 Token Choice 路由器，
使得最终的路由器在保持 Token Choice 语义的同时
具有更均匀的负载分配特性。

\subsection{其他路由策略}

\textbf{Hash 路由}是最简单的基线：
根据 token 的哈希值确定性地分配到专家。
不需要可学习的路由器，负载完全均衡，
但缺乏语义感知，
代码 token 和诗歌 token 可能被分到同一个专家。
实验表明 Hash 路由的性能显著低于学习型路由。

\textbf{随机路由}也是一个有用的基线：
随机均匀地将 token 分配到专家。
在某些配置下，随机路由的表现出人意料地不差，
因为即使路由是随机的，
专家本身仍然可以通过梯度下降
学到对不同类型输入的通用处理能力。
这提示了一个有趣的问题：
路由器到底有多重要？
答案是：在大规模训练中，学习型路由\textbf{显著}优于随机路由，
但差距不如你直觉认为的那么大。

这引出了一个关于 MoE 的深层问题：
\textbf{路由的质量和专家的质量，哪个更重要？}
Switch Transformer 的消融实验表明，
当专家数量很多（$>$64）时，
路由质量是决定性因素：
因为每个专家只需要处理总 token 数的一小部分，
它们有足够的参数去学好自己那一份工作。
但当专家数量很少（$<$8）时，
每个专家承担的负载很重，
专家本身的容量成为瓶颈，
路由的改进带来的收益相对有限。
这个规律对架构设计有直接指导意义：
如果你选择了大量细粒度专家的路线（如 DeepSeek 的 256 专家），
路由器的精度至关重要；
如果选择了少量大专家的路线（如 Mixtral 的 8 专家），
每个专家的内部容量更为关键。

\textbf{Soft MoE}（Puigcerver 等人，2023）% NEW REF: puigcerver2024softmoe = Puigcerver et al., "From Sparse to Soft Mixtures of Experts", arXiv:2308.00951, 2023
走了一条完全不同的路线：
取消离散的 Top-K 选择，
让每个 token 被\textbf{所有}专家以不同权重处理。
具体来说，Soft MoE 将一组 token 线性组合后
送入每个专家，而不是将单个 token 路由到特定专家。
这消除了离散路由带来的梯度不连续性，
也天然避免了负载不均衡，
每个专家处理的是 token 的加权混合，而非原始 token。

Soft MoE 的问题在于它破坏了 MoE 的稀疏性优势：
如果每个 token 对所有专家都有非零贡献，
那么所有专家都必须参与计算，
计算量回到了稠密模型的水平。
在实践中，Soft MoE 通常结合注意力分数的稀疏化
来近似恢复部分计算节省，
但其推理效率仍远不如 Top-K 路由。
Google 的 ViT-MoE 在视觉任务上使用了 Soft MoE
并取得了不错的效果~\cite{puigcerver2024softmoe}，
但在语言模型领域，Token Choice 的 Top-K 路由
因其清晰的稀疏性优势仍然占主导。

\subsection{动态路由与可微路由（2025--2026 前沿）}
\label{sec:advanced-routing}

标准的 Top-K 路由有一个根本局限：
\textbf{每个 token 固定激活 $K$ 个专家}，
无论 token 的复杂度如何。
简单的 token（如 “the”）可能只需要 1 个专家就够了，
而复杂的 token（如数学公式中的关键符号）
可能需要 5 个以上的专家协同处理。

\textbf{DynaMoE}~(2026) 放松了这两个约束：% NEW REF: gulme2026dynamoe = DynaMoE: Dynamic Token-Level Expert Activation with Layer-Wise Adaptive Capacity, Guelmez, arXiv 2026
（1）每个 token 激活的专家数量根据输入复杂度
动态变化（从最少 1 个到层级相关的最大值）；
（2）不同层的专家总容量可以不同，
底层可能需要更多专家处理低级特征的多样性，
顶层可能需要更少但更专精的专家。
DynaMoE 提出了 6 种容量调度策略
（线性递增、线性递减、中间峰值等），
让模型自动学习最优的层间容量分配。

\textbf{DirMoE} 是 ICLR 2026 上提出的% NEW REF: dirmoe2026 = DirMoE: Differentiable Probabilistic Router for Mixture-of-Experts, ICLR 2026
\textbf{完全可微}的概率路由器。
传统的 Top-K 路由不可微，
“选择第 3 和第 7 个专家” 这个离散决策
无法产生干净的梯度信号。
DirMoE 将路由分解为两个可微的概率过程：
\begin{itemize}[noitemsep]
  \item \textbf{Bernoulli 过程}决定每个专家是否被激活
        （连续松弛的 Gumbel-Sigmoid）。
  \item \textbf{Dirichlet 过程}决定被激活专家之间的权重分配。
\end{itemize}
这种分离使得路由在端到端训练中拥有干净的梯度流。
DirMoE 还引入了基于 Simpson 指数的“稀疏度旋钮”，
通过一个标量超参数精确控制期望的活跃专家数量，
\textbf{无需辅助损失}来维持负载均衡。
实验表明 DirMoE 在匹配标准 MoE 吞吐量的同时
实现了更稳定的路由和更均匀的专家利用率。

\textbf{L2R}（Low-Rank \& Lipschitz-Controlled Routing）% NEW REF: l2r2026 = L2R: Low-Rank and Lipschitz-Controlled Routing for Mixture-of-Experts, Yang et al., arXiv 2026
从另一个角度改进路由：
在低秩潜在空间中执行路由决策，
并通过 Lipschitz 控制确保路由函数的平滑性。
直觉上，标准的线性路由器在高维空间中
可能因“角度集中”现象而产生
几乎相同的分数，微小的输入变化导致路由决策翻转。
L2R 通过约束路由函数的 Lipschitz 常数，
确保了路由决策的稳定性。
具体来说，Lipschitz 约束意味着
$\|f(\mathbf{x}_1) - f(\mathbf{x}_2)\| \leq L \cdot \|\mathbf{x}_1 - \mathbf{x}_2\|$，
其中 $L$ 是 Lipschitz 常数。
较小的 $L$ 意味着路由函数对输入扰动更加“平滑”，
两个相似的 token 会被路由到相同的专家集合，
而不是因为微小的数值差异被送到完全不同的专家。
这种稳定性在训练中尤为重要：
它防止了路由器在梯度更新后
对大量 token 做出不同的路由决策（路由振荡），
从而提升了训练的整体稳定性。

L2R 的低秩投影也带来了一个额外好处：
降低路由器的参数量和计算量。
对于 256 个专家、$d = 8192$ 的配置，
标准路由器的参数量为 $8192 \times 256 \approx 2M$。
L2R 通过 $r = 128$ 的低秩分解，
将参数量降低到 $8192 \times 128 + 128 \times 256 \approx 1.1M$，
节省约 45\%。
虽然路由器的参数量相对于专家参数本身微不足道，
但在推理的 decode 阶段
（每步需要执行一次路由决策），
更小的路由器意味着更低的延迟。

这些前沿的路由改进解决的是同一个根本问题：
\textbf{如何让路由决策与 token 的实际需求匹配，
同时保持训练的稳定性和可微性}。
传统的 Top-K 路由是一个粗暴的近似：
用一个固定的 $K$ 来处理所有 token，
用不可微的 argmax 来做离散选择，
再用辅助损失来修补负载不均衡。
DynaMoE 解决了固定 $K$ 的问题，
DirMoE 解决了不可微性的问题，
L2R 解决了路由稳定性的问题。
将这三种思想组合，
在可微的概率框架中做动态数量的专家选择，
同时通过 Lipschitz 控制保证稳定性，
是一个自然的但尚未实现的研究方向。

值得注意的是，这些改进的实际收益
在不同模型规模上可能差异很大。
对于 8 个专家的 Mixtral 量级模型，
DynaMoE 的动态 $K$ 带来的改进可能有限，
因为当 $K$ 在 1--4 之间变化时，
组合空间本来就不大。
但对于 256 个以上专家的模型，
动态 $K$ 的价值更加显著，
让简单 token 只激活 2--3 个专家（节省大量计算），
让复杂 token 激活 12--15 个专家（提升质量），
这种自适应性在大规模配置下收益更大。

用代码来看，MoE 路由器的核心逻辑可以用几行实现：

\begin{lstlisting}[caption={MoE router implementation}]
class MoERouter(nn.Module):
    def __init__(self, d_model, n_experts, top_k):
        super().__init__()
        self.gate = nn.Linear(d_model, n_experts, bias=False)
        self.top_k = top_k

    def forward(self, x):
        # x: (batch, seq_len, d_model)
        scores = self.gate(x)  # (batch, seq_len, n_experts)

        # Select top-k experts per token
        topk_scores, topk_ids = scores.topk(self.top_k, dim=-1)

        # Normalize scores (softmax over selected experts only)
        topk_weights = F.softmax(topk_scores, dim=-1)

        return topk_weights, topk_ids
\end{lstlisting}

这段代码展示了路由的三个步骤：
线性投影得到每个专家的分数，
TopK 选出得分最高的 $K$ 个专家，
在选出的专家上做 softmax 得到归一化的权重。
这些权重决定了各专家输出的贡献比例。
其中 $g_i$ 是归一化后的路由权重，
只在被选中的 TopK 专家上计算 softmax。

\section{负载均衡：MoE 训练的核心挑战}

\subsection{为什么负载不均衡是灾难性的}

MoE 训练中最棘手的问题是\textbf{专家坍缩}
（expert collapse）：
路由器学会将绝大多数 token 发送到少数几个「强」专家，
而大部分专家几乎不接收任何 token，逐渐退化为随机噪声。

为什么这是灾难性的？两个层面：

\textbf{计算浪费}：
假设有 64 个专家分布在 64 块 GPU 上，
如果 80\% 的 token 都被路由到 4 个专家，
那么只有 4 块 GPU 在满负荷工作，
另外 60 块 GPU 基本在空转，
你花了 64 块 GPU 的钱，只得到了 4 块 GPU 的效果。
在分布式训练中，整个系统的速度受最慢节点（最忙的专家）制约，
所以负载不均衡直接导致训练吞吐量暴跌。

\textbf{知识容量浪费}：
MoE 的核心优势是「用总参数量换知识容量」。
如果大部分专家闲置，等价于这些参数不存在，
671B 参数的 MoE 模型可能退化为
只有几十 B 参数的效果。

\subsection{传统方案：辅助损失}

最早的解决方案~\cite{shazeer2017moe}
是在训练目标中加入一个辅助损失来惩罚不均衡：
\begin{equation}
  L_{\mathrm{total}} = L_{\mathrm{LM}} + \alpha \cdot L_{\mathrm{balance}}
  \label{eq:aux-loss}
\end{equation}
其中辅助损失的形式为：
\begin{equation}
  L_{\mathrm{balance}} = N \cdot \sum_{i=1}^{N} f_i \cdot P_i
  \label{eq:balance-loss}
\end{equation}
这里 $f_i$ 是实际被分配到专家 $i$ 的 token 比例，
$P_i$ 是所有 token 对专家 $i$ 的平均路由概率。
$N$ 是专家数量，用于归一化。

直觉上，$f_i \cdot P_i$ 的乘积形式
鼓励两者同时尽可能均匀：
当所有专家均匀接收 token 时
（$f_i = P_i = 1/N$），
$\sum f_i P_i = N \times (1/N)^2 = 1/N$，达到最小值。
当负载高度不均衡时（所有 token 集中在一个专家），
$\sum f_i P_i = 1$，达到最大值。

Switch Transformer~\cite{fedus2022switch} 简化了这个公式，
使用 $\alpha = 0.01$ 作为默认系数，
并证明即使是 top-1 路由，
只要辅助损失设置得当，训练也能保持稳定。

除了 $f_i \cdot P_i$ 的“重要性--负载”乘积形式，
还有几种变体被广泛使用。
\textbf{Z-loss}~\cite{fedus2022switch} 惩罚路由 logits 的绝对值：
$L_z = \frac{1}{B} \sum_{b=1}^B (\log \sum_i e^{s_i^{(b)}})^2$。
它的目的不是直接促进负载均衡，
而是防止路由 logits 变得过大，
过大的 logits 会导致 softmax 输出趋向 one-hot，
使得路由器“过于自信”，
拒绝探索其他专家的可能性。
Z-loss 可以与标准辅助损失联合使用，
$\alpha$ 通常设为 0.001 量级。

\textbf{Importance loss} 只考虑路由概率的均匀性：
$L_{\text{imp}} = \text{CV}(P_1, P_2, \ldots, P_N)^2$，
其中 CV 是变异系数（标准差除以均值）。
它不直接约束实际负载，
而是确保路由器的“意图”是均匀的。
实际负载可能因为 Top-K 的离散化而与路由概率不一致，
但长期来看，路由概率的均匀性
会引导实际负载趋向均匀。

但辅助损失有一个\textbf{根本矛盾}：
$\alpha$ 太大，负载均衡得到满足，
但辅助损失的梯度「污染」了语言建模的梯度，
路由器不再纯粹地为质量最优而路由，
而是在「质量」和「均衡」之间妥协，模型质量下降。
$\alpha$ 太小，辅助损失形同虚设，
专家坍缩照常发生。
在实践中，$\alpha$ 的调节需要大量实验，
且没有一个值能适用于训练的所有阶段。

\subsection{DeepSeek 的无辅助损失方案}

DeepSeek-V3~\cite{deepseekv3} 提出了一个优雅的替代方案：
\textbf{不加辅助损失}，
而是给每个专家的路由分数加一个\textbf{不参与梯度计算}的偏置项 $b_i$：
\begin{equation}
  s_i = W_{\mathrm{router}} \mathbf{x} + b_i
  \quad (\text{$b_i$ 不参与反向传播})
\end{equation}
偏置的更新规则极其简单：
如果第 $i$ 个专家过载，$b_i \leftarrow b_i - \gamma$；
如果空闲，$b_i \leftarrow b_i + \gamma$（$\gamma \approx 10^{-3}$）。

因为 $b_i$ 不在计算图中，
语言建模的梯度完全不受干扰，
训练目标纯粹是预测下一个 token，没有任何妥协。
2025 年芝加哥大学的理论论文证明了这种方法等价于
原始--对偶分配问题的在线求解，具有单调改进保证。

这个方法的优雅之处在于它将两个不同的优化问题\textbf{解耦}了：
语言建模（预测下一个 token）通过标准的梯度下降优化，
负载均衡通过偏置项的启发式规则调整。
两者互不干扰，语言建模的梯度不受负载均衡的影响，
负载均衡的调整也不需要通过梯度传播。
这就像一个工厂的生产线和排班表是分开管理的：
生产线的优化不需要考虑排班约束，
排班表的调整也不需要改变生产流程。

2025 年芝加哥大学的理论分析~\cite{deepseekv3}
进一步揭示了为什么偏置调节在实践中如此有效。
论文将 DeepSeek 的偏置更新规则
映射到一个\textbf{在线线性规划}的对偶形式中，
证明了以下性质：
（1）偏置值 $b_i$ 随训练步数单调收敛到一个稳定点；
（2）在收敛点处，所有专家的负载与最优均衡之间的偏差
以 $O(1/\sqrt{T})$ 的速率衰减（$T$ 为训练步数）；
（3）该收敛不依赖于语言建模损失的具体形式，
即使模型在不同训练阶段学习不同类型的知识
（先学语法，再学事实，再学推理），
偏置调节都能自适应地跟踪变化的路由需求。
这些理论保证解释了 DeepSeek-V3
在 14.8T token 的漫长训练过程中
偏置方案始终保持稳定的原因。

\begin{keybox}[负载均衡方法对比]
\begin{center}\small
\renewcommand{\arraystretch}{1.3}
\begin{tabular}{lccl}
  \toprule
  \rowcolor{figLightGray}
  \textbf{方法} & \textbf{需要辅助损失} & \textbf{影响 LM 梯度}
  & \textbf{代表模型} \\
  \midrule
  辅助损失 & 是 & 是（$\alpha$ 敏感） & Switch Transformer \\
  Expert Choice & 否 & 部分 & Expert Choice Routing \\
  偏置项调节 & 否 & 否 & DeepSeek-V3 \\
  \bottomrule
\end{tabular}
\end{center}
\end{keybox}

\section{MoE 的训练动力学}

理解 MoE 如何在训练中“学会分工”
是理解整个架构的关键。
路由器和专家之间存在一种微妙的协同进化关系：
路由器的决策影响每个专家收到的数据分布，
而专家的性能反过来影响路由器的梯度更新。
这种耦合使得 MoE 的训练比稠密模型更加复杂和不稳定。

\subsection{专家专业化的涌现}

MoE 训练中最令人着迷的现象之一
是\textbf{专家自发地发展出不同的专长}。
没有任何人为标注告诉路由器
“把代码 token 送到专家 17，把数学 token 送到专家 42”，
这种分工完全是训练过程中涌现出来的。

多项研究对训练完成的 MoE 模型的路由模式
进行了系统性的分析，揭示了几个有趣的规律。

\textbf{语言专业化}。
在多语言训练的 MoE 模型中，
不同专家经常发展出对不同语言的偏好。
例如，某些专家可能主要处理中文 token，
另一些专家偏好英文或法文。
但这种语言专业化并非绝对，
同一个专家可能同时擅长中文和日文
（它们共享大量汉字），
而一些“通用”专家在所有语言上都有中等活跃度。
这种模式暗示路由器学到了某种
“语言家族”的隐式表示。

\textbf{任务专业化}。
某些专家在处理代码、数学公式、
或特定格式（如 JSON、HTML）的文本时
被显著更频繁地激活。
在 DeepSeek-V3 的分析中~\cite{deepseekv3}，
研究者发现某些专家几乎只在
Python 代码的特定模式（如列表推导、装饰器）中被激活，
而在自然语言文本中完全沉默。
这种极端的专业化表明，
细粒度专家（256 个小专家 vs 8 个大专家）
确实能带来更精准的分工。

\textbf{层间专业化差异}。
底层的专家倾向于学习表面特征（token 类型、标点、格式），
中间层的专家开始分化出语义领域（科学、法律、代码），
顶层的专家则可能发展出功能性分工
（检索事实、执行推理、生成语法正确的输出）。
这种层间差异暗示了不同层对“专业化”的定义不同，
底层的“专业”是语法层面的，
顶层的“专业”是语义和功能层面的。

\subsection{“强者愈强”问题}

MoE 训练初期存在一个危险的正反馈循环：
某个专家因为随机初始化的微小优势
在某些 token 上表现稍好，
路由器于是给它分配更多这类 token，
它在这类 token 上获得更多训练数据而进一步改善，
路由器就给它分配更多更多的 token……
最终，少数专家垄断了大部分流量，
其余专家沦为“僵尸专家”，
参数存在但几乎不参与有效计算。

这个“强者愈强”（rich-get-richer）动态
在经济学中被称为马太效应，
在 MoE 训练中表现得尤为剧烈。
问题的根源在于路由器的梯度信号：
只有被选中的专家能收到梯度更新，
未被选中的专家的梯度为零。
如果一个专家在训练早期连续多个 batch 都未被选中，
它的参数保持在随机初始化状态，
路由器对它的评分自然更低，
形成了一个越来越难以打破的恶性循环。

\textbf{容量因子}（capacity factor）是应对这一问题的早期方案。
每个专家设置一个最大处理容量 $C = c \cdot (B / N)$，
其中 $B$ 是 batch 中的总 token 数，
$N$ 是专家数量，$c > 1$ 是容量因子。
当某个专家接收的 token 数超过 $C$ 时，
多余的 token 被“溢出”（overflow），
直接跳过该 MoE 层通过残差连接传递。
这机制防止了少数专家处理过多 token 的情况，
但被溢出的 token 在该层得不到有效处理，
质量可能受损。

Switch Transformer~\cite{fedus2022switch} 通过大量实验
确定了 $c = 1.25$（即允许每个专家超额 25\%）
在训练质量和负载均衡之间取得了良好平衡。
更高的容量因子意味着更多 token 能被有效处理，
但也意味着更大的内存占用（每个专家需要预留更多缓冲区）。
DeepSeek-V3 通过引入无辅助损失的偏置调节方案，
在很大程度上避免了对容量因子的依赖。

\subsection{训练不稳定性：Loss Spike 问题}

大规模 MoE 训练中一个令人头疼的问题是
\textbf{loss spike}，
训练损失突然剧烈增大，
然后经过数百甚至数千步才恢复到正常水平。
loss spike 在稠密模型训练中也偶尔发生，
但在 MoE 中更加频繁和严重~\cite{shazeer2017moe}。

MoE 的 loss spike 有一个独特的来源：
\textbf{路由振荡}（routing oscillation）。
当路由器的权重发生较大更新时，
大量 token 的路由决策可能同时改变，
原本发给专家 A 的 token 突然被重定向到专家 B。
专家 B 接收到一批它“不习惯”的数据，
输出质量骤降；专家 A 则突然失去了主要数据来源，
梯度信号变弱。
这种全局性的路由重组就像一家公司
突然把所有员工调换岗位，
每个人都需要时间适应新角色，
整体效率暂时暴跌。

Kimi K2 在训练万亿参数 MoE 时~\cite{kimiTeam2025kimik2}，
通过改进的 MuonClip 优化器实现了“零 loss spike”，
这在万亿参数规模上是一个重要的里程碑。
MuonClip 的核心思想是
对路由器权重的更新步长施加更严格的 Lipschitz 约束，
防止路由决策在单步更新中发生过大变化。
直觉上，这就像给路由器戴上了“手铐”，
它可以逐步调整分配策略，
但不允许在一步之内做出剧烈的重新分配。

\subsection{MoE 微调的挑战}

MoE 模型的微调（fine-tuning）面临独特的困难，
这是阻碍 MoE 在下游应用中大规模采纳的重要因素之一。

\textbf{LoRA 与 MoE 的不兼容性}。
LoRA（Low-Rank Adaptation）是目前最流行的参数高效微调方法，
但它在 MoE 上的应用面临一个设计抉择：
应该给哪些部分加 LoRA 适配器？

\begin{itemize}[noitemsep]
  \item \textbf{方案 A：只给路由器加 LoRA}。
        参数量极小（路由器本身就很轻量），
        但仅改变路由策略，
        不改变专家的处理能力，
        相当于重新排班但不培训员工。
  \item \textbf{方案 B：给每个专家都加 LoRA}。
        效果最好，但 256 个专家 $\times$ 每个专家的 LoRA 适配器
        = 大量的额外参数。
        而且由于不同专家处理不同类型的数据，
        每个专家的 LoRA 需要独立训练，
        训练数据被分散到多个专家上，
        每个专家的 LoRA 只看到总数据的一小部分。
  \item \textbf{方案 C：只给共享专家和高频路由专家加 LoRA}。
        这是一种实用的折中：
        共享专家处理所有 token，
        它的 LoRA 能影响所有输出；
        高频路由专家（被最常激活的前 $k$ 个）
        覆盖大部分 token。
        实验表明，给共享专家 + top-20\% 的路由专家加 LoRA
        能恢复全参数微调约 90\% 的效果，
        同时参数量不到全参数的 3\%。
\end{itemize}

\textbf{路由漂移}（routing drift）是 MoE 微调中的另一个问题。
预训练时学到的路由模式
在微调数据分布发生变化时可能失效。
例如，一个在通用语料上预训练的 MoE 模型
被微调为专注于医学文本，
原来处理代码和数学的专家可能在微调后变得几乎不被激活，
而处理自然科学文本的专家被过度激活。
这种路由模式的剧变可能导致部分专家“退化”，
微调后的模型在非医学任务上性能严重下降
（灾难性遗忘），
甚至在医学任务上也不如预期
（因为路由器需要时间适应新的分配模式）。

应对路由漂移的方法包括：
\textbf{冻结路由器}（只微调专家权重，保持路由不变）、
\textbf{渐进式路由适应}（用很小的学习率更新路由器，
大学习率更新专家）、
以及 \textbf{路由正则化}（在微调损失中加入惩罚项，
限制路由分布与预训练时的偏差）。
目前没有一种方法在所有场景下都最优，
MoE 微调仍然是一个活跃的研究方向。

\section{DeepSeek 的三大创新}

DeepSeek-V3~\cite{deepseekv3} 在 MoE 上做了三项改变游戏规则的创新。

\subsection{细粒度专家}

传统 MoE 使用少量大专家（如 Mixtral 的 8 个专家，每个约 7B 参数）。
DeepSeek 使用 \textbf{256 个小专家}（每个仅约 0.5B 参数）
加 1 个\textbf{共享专家}（处理所有 token 的通用知识）：
\begin{equation}
  \mathbf{y} = \mathrm{SharedExpert}(\mathbf{x})
  + \sum_{i \in \mathrm{Top8}} g_i \cdot \mathrm{Expert}_i(\mathbf{x})
\end{equation}

256 个专家中选 8 个，组合数为
$\binom{256}{8} \approx 4.4 \times 10^{9}$，
相比 Mixtral 的 $\binom{8}{2} = 28$ 种组合，
组合空间扩大了 $10^{8}$ 倍。
这意味着模型可以为不同类型的输入激活完全不同的专家组合，
实现更精细的专业化。

为什么细粒度专家更好？三个原因：

\textbf{更灵活的组合}。
8 个大专家中选 2 个，
就像只有 8 门课的课程表中选 2 门，
选择空间很有限。
256 个小专家中选 8 个，
就像 256 门课中选 8 门，
可以精确匹配每个学生（token）的需求。

\textbf{更好的专业化}。
小专家被迫在一个更窄的领域内做到精通，
因为它们的参数量有限，无法成为「通才」。
大专家则可能在多个不相关的领域之间分散注意力。
实验表明，细粒度专家的专业化程度更高，
某些专家可能只处理 Python 代码中的列表推导，
另一些只处理数学证明中的符号推理。

\textbf{相似的计算成本}。
关键洞察：8 个小专家的总计算量
$\approx$ 2 个大专家的计算量（如果大小比例为 1:4）。
所以细粒度专家不增加每个 token 的计算成本，
它只是用更多的小模块替代了少数大模块。

细粒度专家的设计还有一个不太明显但重要的好处：
\textbf{它改善了 GPU 上的计算效率}。
大专家的单次前向传播涉及大型矩阵乘法（如 7B 参数的 FFN），
而小专家的矩阵乘法规模更小。
这看起来似乎不利于 GPU 利用率
（GPU 更擅长大矩阵运算），
但实际情况更加微妙。
在 batched inference 中，
多个 token 可能被路由到同一个小专家，
它们的计算可以合并为一个 batch，
如果一个小专家收到 $B$ 个 token 的 batch，
有效矩阵大小是 $[B, d_{\text{expert}}] \times [d_{\text{expert}}, d_{\text{ff}}]$，
其中 $B$ 可以足够大以饱和 GPU 的计算能力。
关键在于 $B$ 的大小取决于有多少 token 被路由到该专家，
在 token-level batching 的推理系统中（如 vLLM），
来自不同请求的 token 可以被打包在一起，
使得即使是小专家也能获得足够大的 batch。

\subsection{共享专家}

在上面的公式中，
$\mathrm{SharedExpert}(\mathbf{x})$ 是一个始终被激活的专家：
不管 token 是什么类型，它都参与计算。

为什么需要共享专家？
在纯路由的 MoE 中，
每个路由专家都需要独立学会一些「基础知识」
（如语法规则、常见词的语义等），
因为不能保证某个 token 每次都被路由到同一个专家。
这导致了跨专家的冗余，
256 个专家中可能有 200 个都各自学了一遍「the 后面通常跟名词」。

共享专家将这些通用知识集中到一个始终激活的模块中，
路由专家只需学习\textbf{差异化}的专业知识。
这既减少了冗余，又让路由专家能更专注于自己的领域。

从残差连接的视角看，
共享专家提供了一个强大的「默认输出」，
即使路由器把 token 送到了不太合适的专家，
共享专家的输出仍然保证了基本的质量。

\subsection{MLA + MoE：注意力与专家的协同设计}

DeepSeek-V3 的另一个关键设计决策是
将 \textbf{MLA}（Multi-head Latent Attention）
与 MoE 结合使用，
这在 KV cache 压缩和稀疏计算之间实现了双重效率提升。

MLA 的核心思想是在注意力计算中引入潜在空间压缩。
传统的多头注意力（MHA）为每个头独立存储 K 和 V 向量，
GQA（Grouped Query Attention）通过让多个 query 头共享
一组 KV 头来减少 KV cache。
MLA 走得更远，
它将所有 KV 头的信息压缩到一个低维的潜在向量
$\mathbf{c}_{\mathrm{KV}} \in \mathbb{R}^{d_c}$（$d_c \ll n_h \cdot d_h$），
在推理时只需缓存这个低维向量，
需要时再通过线性投影恢复出各头的 K 和 V。

MLA 与 MoE 的组合产生了一种“双重稀疏”架构：
\begin{itemize}[noitemsep]
  \item \textbf{注意力层}：MLA 通过潜在空间压缩
        将 KV cache 大小降低了约 $4\times$--$8\times$。
        对于 128K 上下文，这意味着
        KV cache 从几十 GB 降到几 GB。
  \item \textbf{FFN 层}（被 MoE 替换）：
        只有 Top-K 个专家参与计算，
        计算量降低约 $N/K$ 倍。
\end{itemize}

两种稀疏性作用于 Transformer 的不同组件
（注意力 vs FFN），互不干扰，效果叠加。
这解释了 DeepSeek-V3 为什么能以
$\sim$\$5.6M 的成本训练出 671B 参数的模型，
不仅因为 MoE 减少了前向计算，
还因为 MLA 减少了注意力计算中的内存和带宽需求。

从 DeepSeek-V2 到 V3 的演进清晰地展示了
MLA + MoE 的协同优化路径：
V2 首次将 MLA 与 MoE 结合，
验证了技术可行性但在训练稳定性上遇到了挑战。
V3 引入了无辅助损失的负载均衡、FP8 训练、
以及 DualPipe 计算--通信重叠等工程创新，
解决了大规模部署的稳定性和效率问题。
这种迭代路径表明，
\textbf{架构创新和工程优化必须同步推进}，
一个好的架构设计如果没有配套的训练基础设施，
很难在实践中发挥全部潜力。

\subsection{DeepEP：专用通信库}

MoE 的一个工程瓶颈是 All-to-All 通信：
每个 token 需要被发送到对应专家所在的 GPU，
计算完成后再发送回来。
这种细粒度的跨 GPU 通信模式不同于数据并行的 All-Reduce，
标准通信库的效率很低。

为什么 All-to-All 比 AllReduce 更难优化？
AllReduce 中每块 GPU 发送和接收的数据量是确定性的：
总是 $M/N$ 的数据。
All-to-All 中数据量取决于路由决策，
如果某个专家特别热门，
对应 GPU 可能收到 10 倍于平均水平的数据，
而其他 GPU 几乎收不到数据。
这种\textbf{不均匀的通信模式}
让带宽利用和延迟优化变得困难。

DeepSeek 开源的 DeepEP 提供两种模式：

\textbf{Normal 模式}用于训练和 prefill（高吞吐量优先）：
将大量 token 的路由请求打包成批次，
利用 NVLink（节点内）和 InfiniBand（节点间）的
分层通信策略，
同一节点内的专家通信走 NVLink（900\,GB/s），
跨节点的专家通信走 IB（400\,Gb/s），
两者可以并行进行。

\textbf{Low-Latency 模式}用于推理 decode
（通过 hook-based 通信--计算重叠实现零额外 SM 占用）。
在 decode 阶段，每个时间步只有 1 个 token 需要路由，
通信量很小但对延迟极其敏感。
DeepEP 利用 GPU 的硬件 hook 机制
将通信操作卸载到 NVLink/IB 控制器，
SM 完全不参与通信，
计算和通信真正同时进行，零开销。

DeepEP 的开源（Apache 2.0 许可）
对整个 MoE 生态系统产生了显著影响。
在此之前，高效的 MoE All-to-All 通信
基本上是 NVIDIA NCCL 库的“黑箱”，
用户只能调整几个有限的参数。
DeepEP 的开源让研究者可以深入了解
MoE 通信的瓶颈所在，并进行针对性的优化。
多个后续项目（包括 Megatron Core MoE~\cite{yan2026megatronmoe}）
在其 All-to-All 通信实现中
借鉴了 DeepEP 的设计思想。

值得提及的是 DeepSeek 的\textbf{DualPipe} 优化，
它与 DeepEP 配合工作。
DualPipe 将流水线并行的前向和后向传播
重叠在时间维度上，
当一组 GPU 在做第 $t$ 步的后向传播时，
另一组 GPU 同时开始第 $t+1$ 步的前向传播。
在 MoE 场景下，DualPipe 与 DeepEP 的组合
实现了“计算--通信--内存”三者的同时优化：
计算通过 DualPipe 的流水线重叠，
通信通过 DeepEP 的硬件 hook 卸载，
内存通过选择性激活重算来节省。
这种全栈协同优化是 DeepSeek-V3
能以极低成本完成训练的关键工程因素之一。

\section{MoE 在实践中}

\subsection{Mixtral 8x7B：MoE 的开源先驱}

Mixtral 8x7B~\cite{jiang2024mixtral} 是 Mistral AI 发布的
第一个广泛使用的开源 MoE 模型。
它的规格说明了 MoE 的效率优势：
\begin{itemize}[noitemsep]
  \item \textbf{总参数}：47B（8 个 FFN 专家 $\times$ $\sim$5.5B + 非 FFN 参数）
  \item \textbf{活跃参数}：13B（每个 token 激活 2 个专家）
  \item \textbf{性能}：在大部分基准测试上
        匹配或超越 LLaMA 2 70B（稠密模型）
  \item \textbf{推理成本}：约为 LLaMA 2 70B 的 $1/5$
\end{itemize}

47B 总参数 vs 70B 稠密参数，
Mixtral 的参数总量反而更少，
但它用 MoE 的稀疏激活实现了类似甚至更好的性能，
同时推理时只需要 13B 的计算量。

Mixtral 的成功对开源社区有深远的影响。
在 Mixtral 之前，开源模型与闭源模型
（GPT-4、Claude）之间存在巨大的性能差距。
Mixtral 证明了 MoE 可以让较小的开源模型
在推理效率上接近甚至超越更大的稠密模型，
这给资源有限的研究团队和创业公司
提供了一条可行的追赶路径。

从工程实践角度看，Mixtral 也揭示了 MoE 部署的挑战。
虽然推理时只有 13B 参数是活跃的，
但 47B 的总参数仍然需要全部加载到内存中。
在 8-bit 量化下，Mixtral 需要约 50\,GB 内存，
可以装入单块 H100（80\,GB），
但在消费级 GPU（如 RTX 4090，24\,GB）上
则必须使用 4-bit 量化或部分卸载到 CPU。
Mixtral 的这个经验教训，
“MoE 减少计算但不减少内存”，
成为后续所有 MoE 模型部署规划的核心考量。

\textbf{Mixtral 8x22B} 是 Mistral AI 的后续更新，
总参数扩大到 141B，活跃参数 39B，
在多项基准测试上进一步缩小了与 GPT-4 的差距。
这种从 8x7B 到 8x22B 的扩展路径展示了 MoE 的优雅之处：
不改变专家数量和路由策略（仍然是 8 个专家、Top-2），
仅通过增大每个专家的参数量
就实现了性能的大幅提升，
就像不改变公司的部门结构，
只给每个部门配备更多更强的员工。

\subsection{DeepSeek-V3：MoE 的巅峰}

DeepSeek-V3~\cite{deepseekv3} 将 MoE 推向了新的高度：
\begin{itemize}[noitemsep]
  \item \textbf{总参数}：671B
  \item \textbf{活跃参数}：37B（256 路由专家选 8 + 1 共享专家）
  \item \textbf{训练数据}：14.8T token
  \item \textbf{训练成本}：$\sim$\$5.6M（2048 块 H800，约 2 个月）
  \item \textbf{性能}：在多项基准测试上达到或超越 GPT-4o、Claude 3.5 Sonnet
\end{itemize}

$\$$5.6M 的训练成本令整个行业震惊，
它比同等性能的 LLaMA 3 405B 的估计训练成本低约 20 倍。
成本优势来自三方面叠加：
MoE 架构（37B 活跃 vs 405B 全激活）、
FP8 训练（减少 $\sim$50\% 内存和 $\sim$30\% 计算）、
以及精细的工程优化（DualPipe、DeepEP 等）。

DeepSeek-V3 的成功引发了一个更深层的行业讨论：
\textbf{MoE 是否改变了 AI 训练的经济学？}
传统上，训练一个 frontier 模型需要数亿美元的投入，
GPT-4 的训练成本估计在 \$100M 以上，
Google Gemini 1.5 的估计成本也在类似量级。
DeepSeek-V3 以 \$5.6M 达到了接近的性能水平，
这意味着 frontier AI 的门槛可能比之前认为的低得多。

这对行业格局有深远影响。
如果训练成本从数亿美元降到数百万美元，
那么更多的机构（大学、中型公司、政府研究机构）
都有能力训练自己的 frontier 模型。
MoE 架构在其中扮演了关键角色：
它让模型可以在较少的 GPU 上训练更长时间（因为活跃参数少），
而不是必须用大规模集群做短期的全参数训练。

当然，\$5.6M 的数字需要审慎理解。
这只包括最终训练的计算成本，
不包括 DeepSeek 在 V1、V2 阶段积累的
大量架构搜索和超参数调优成本，
也不包括团队的人力成本和数据处理成本。
但即便如此，MoE 对训练效率的贡献是不可否认的。

\subsection{LLaMA 4：Meta 的 MoE 转型}
\label{sec:llama4-moe}

2025 年 4 月，Meta 发布了 LLaMA 4 系列，% NEW REF: llama4 = The Llama 4 Herd, Meta AI, 2025
这是 LLaMA 家族首次采用 MoE 架构，
标志着 Meta 从稠密模型全面转向 MoE。

LLaMA 4 包含三个变体：
\begin{center}\small
\renewcommand{\arraystretch}{1.3}
\begin{tabular}{lcccl}
  \toprule
  \rowcolor{figLightGray}
  \textbf{模型} & \textbf{总参数} & \textbf{活跃参数}
  & \textbf{专家数 / Top-K} & \textbf{特色} \\
  \midrule
  Scout & 109B & 17B & 16 / 1 & 10M token 上下文 \\
  Maverick & 400B & 17B & 128 / 1 & 原生多模态 \\
  Behemoth & $\sim$2T & 288B & 16 / 2 & 训练中（未发布） \\
  \bottomrule
\end{tabular}
\end{center}

LLaMA 4 的 MoE 设计与 DeepSeek 有显著差异：

\textbf{Top-1 路由}。
Scout 和 Maverick 都使用 Top-1 路由
（每个 token 只发给 1 个专家），
而不是 DeepSeek 的 Top-8。
这是一个大胆的选择：
Top-1 在 Switch Transformer 中曾被证明可行，
但需要更精确的路由器。
Meta 通过从 Behemoth（Top-2）的知识蒸馏
来弥补 Top-1 路由的信息损失。

\textbf{知识蒸馏}。
Scout 和 Maverick 都通过从 Behemoth 蒸馏获得了
显著的质量提升，这是一种“大模型训练小模型”的策略。
Behemoth 本身有 288B 活跃参数（Top-2 路由 16 个专家），
在 STEM 基准测试上超越了 GPT-4.5 和 Claude Sonnet 3.7。

\textbf{超长上下文}。
Scout 支持 \textbf{10M token 上下文窗口}，
目前最长的生产级上下文长度。
这通过高效的 iRoPE 位置编码
和分块注意力机制实现，
但也对 MoE 的路由一致性提出了更高要求。

\subsection{Qwen3-235B：阿里巴巴的 MoE 实践}

Qwen3-235B-A22B 是阿里巴巴 Qwen 团队% NEW REF: qwen3 = Qwen3 Technical Report, Qwen Team, arXiv 2025
发布的旗舰 MoE 模型，采用了与 DeepSeek 相似但不同的设计选择：

\begin{itemize}[noitemsep]
  \item \textbf{总参数}：235B，\textbf{活跃参数}：22B。
  \item \textbf{128 个路由专家，Top-8 路由}，
        与 DeepSeek-V3 的 256 专家 + Top-8 相比，
        专家数减半但每个专家更大。
  \item \textbf{94 层}，使用 GQA（64 头，4 KV 头）。
  \item 支持 128K token 上下文，通过 YaRN 位置编码外推。
\end{itemize}

Qwen3 的一个独特创新是\textbf{思考/非思考双模式}，
模型可以在“思考模式”（生成内部推理链）和
“直接回答模式”之间切换，
用户通过 prompt 控制推理预算。
这对 MoE 路由提出了新要求：
思考模式下的 token 可能需要不同的专家组合
来处理推理链中的中间步骤。

Qwen3-235B 的训练策略也值得关注。
它在第一阶段使用了大规模语料预训练
（与 DeepSeek-V3 的 14.8T token 在同一量级），
然后通过一个四阶段的后训练流程
（SFT $\to$ 推理 RL $\to$ 思考模式融合 $\to$ 通用 RL）
来激活模型的推理能力。
这种“预训练 + 复杂后训练”的范式
暗示了 MoE 模型的质量不仅取决于架构设计，
后训练中路由器的再优化和专家的精调
同样至关重要。
一个有趣的实验发现是：
在 RL 后训练阶段，
部分原本“沉默”的专家被重新激活，
说明推理能力的学习可能需要不同于语言建模的专家组合，
后训练实质上是在\textbf{重新发现}模型中被预训练阶段忽视的专家容量。

\subsection{Kimi K2：万亿参数的 MoE}

Moonshot AI 的 Kimi K2 将 MoE 推到了\textbf{万亿参数}量级：% NEW REF: kimiTeam2025kimik2 = Kimi K2: Open Agentic Intelligence, Kimi Team, arXiv 2025
\begin{itemize}[noitemsep]
  \item \textbf{总参数}：$\sim$1T（1.04T）。
  \item \textbf{活跃参数}：32B。
  \item \textbf{384 个路由专家，Top-8 路由}，
        比 DeepSeek-V3 更多的专家（384 vs 256），
        组合空间为 $\binom{384}{8} \approx 1.8 \times 10^{13}$。
  \item 使用 \textbf{MLA}（Multi-head Latent Attention）
        压缩 KV cache。
  \item 128K token 上下文。
\end{itemize}

Kimi K2 在 MoE 训练上的关键创新是
将 \textbf{Muon 优化器}引入 MoE 训练
（使用改进版 MuonClip），
并在 15.5T token 的全程训练中实现了“零 loss spike”。
这是一个重要的工程里程碑：
万亿参数 MoE 模型的训练稳定性
一直是业界最大的担忧之一。

Kimi K2 还在 MoE 架构上引入了几个值得注意的设计选择。
它使用了与 DeepSeek-V3 相同的 MLA 注意力机制
来压缩 KV cache，
验证了 MLA + MoE 的组合不是 DeepSeek 的“专利”，
而是一种可复现的高效架构范式。
Kimi K2 的 384 个专家数量比 DeepSeek-V3 的 256 个多出 50\%，
这带来了更大的组合空间
但也对路由器的精度提出了更高要求。
为了应对这一挑战，
Kimi K2 在路由器中使用了一个两层的 MLP
（而非标准的单层线性投影），
虽然增加了极少量的参数，
但显著提升了路由决策在 384 个专家中的区分度。

Kimi K2 的另一个创新是\textbf{Agent 能力的原生集成}。
通过在训练数据中大量引入 tool-use 和 multi-step reasoning 数据，
Kimi K2 的路由器学会了为“需要调用工具的 token”
和“需要深度推理的 token”
激活不同的专家组合。
这种“任务感知”的路由行为
在某种程度上模拟了人类大脑中
“执行功能”和“陈述性记忆”使用不同脑区的现象。

\begin{keybox}[2025--2026 主要 MoE 模型对比]
\begin{center}\small
\renewcommand{\arraystretch}{1.3}
\begin{tabular}{lccccc}
  \toprule
  \rowcolor{figLightGray}
  \textbf{模型} & \textbf{总参数} & \textbf{活跃参数}
  & \textbf{专家数} & \textbf{Top-K} & \textbf{共享专家} \\
  \midrule
  Mixtral 8x7B & 47B & 13B & 8 & 2 & 无 \\
  DeepSeek-V3 & 671B & 37B & 256 & 8 & 1 \\
  LLaMA 4 Maverick & 400B & 17B & 128 & 1 & 无$^\dagger$ \\
  LLaMA 4 Behemoth & $\sim$2T & 288B & 16 & 2 & 无$^\dagger$ \\
  Qwen3-235B & 235B & 22B & 128 & 8 & 未公开 \\
  Kimi K2 & 1.04T & 32B & 384 & 8 & 有 \\
  \bottomrule
\end{tabular}
\end{center}

\small $^\dagger$\,LLaMA 4 未明确公开共享专家设计。
\end{keybox}

从表中可以看到一个清晰的趋势：
\textbf{细粒度小专家 + 高 Top-K} 成为主流设计
（DeepSeek、Qwen3、Kimi K2 都使用 Top-8 + 100 个以上的专家）。
LLaMA 4 的 Top-1 路由 + 知识蒸馏是一条独特的路线，
但 Behemoth（未发布的旗舰模型）也回到了 Top-2。

\subsection{架构设计空间的深入对比}

上表只展示了数字，但真正有意义的是这些数字背后的\textbf{设计哲学差异}。
让我们从几个关键维度做更深入的对比。

\textbf{路由策略的分歧}。
行业在路由策略上形成了两个明显的阵营。
“中国阵营”（DeepSeek、Qwen、Kimi）倾向于
大量细粒度专家 + 高 Top-K（8 或更多），
理论基础是更多的组合空间带来更精准的专家匹配。
“Meta 阵营”（LLaMA 4）选择了相反的方向：
较少的专家 + 极低的 Top-K（1 或 2），
辅以大规模蒸馏来弥补信息损失。
两种路线目前没有明确的胜负，
DeepSeek-V3 和 LLaMA 4 在不同基准测试上各有优劣。
但一个有趣的观察是：
Top-1 路由在通信效率上有天然优势，
每个 token 只需发送到一个专家所在的 GPU，
All-to-All 通信量最小化。
这可能是 Meta 选择 Top-1 的工程考量之一，
因为 Meta 的训练集群在跨节点通信带宽上
可能不如 DeepSeek 的定制 InfiniBand 网络。

\textbf{共享专家的角色}。
DeepSeek 和 Kimi K2 都使用了共享专家，
而 Mixtral 和 LLaMA 4 没有。
共享专家的存在改变了路由专家的学习目标，
路由专家只需学习“增量”（差异化知识），
而非“全量”（包括通用知识）。
这使得路由专家可以更小（参数量更低），
因为通用知识的存储被集中到了共享专家。
一个未被广泛注意的细节是：
共享专家的存在实际上改变了 MoE 的有效激活率，
如果共享专家的大小约等于两个路由专家，
那么 DeepSeek-V3 的“有效 Top-K”实际上更接近 10（8+2），
而非名义上的 8。

\textbf{DBRX 的独特设计}。% NEW REF: dbrx2024 = DBRX: A New State-of-the-Art Open LLM, Databricks/Mosaic, 2024
值得一提的是 Databricks 在 2024 年发布的 DBRX，
它使用了 132 个专家和 Top-4 的路由策略，
这是一个介于 Mixtral（少专家 + 低 Top-K）
和 DeepSeek（多专家 + 高 Top-K）之间的折中方案。
DBRX 的 132 个专家数量不是常见的 2 的幂次方，
这暗示 Databricks 可能是通过实验搜索而非理论推导
确定的最优配置。
DBRX 在总参数 132B、活跃参数 36B 的配置下，
性能与当时的 LLaMA 2 70B 和 Mixtral 8x7B 相当或更优，
但由于开源时间较早，
后续被 DeepSeek-V3 和 Qwen3 大幅超越。

\textbf{Grok-1 的另类路线}。
xAI 的 Grok-1 是另一个有趣的数据点：
314B 总参数，8 个专家，Top-2 路由，
设计上更接近 Mixtral 的“少而大”专家风格，
但规模大得多。
Grok-1 的开源代码揭示了一些
不常在论文中讨论的工程细节：
例如它在路由器中使用了额外的 noise injection
来增加探索性（类似 $\epsilon$-greedy 策略），
以及在训练后期逐步降低噪声比例
让路由决策变得更确定。

\section{MoE 推理的工程挑战}
\label{sec:moe-inference}

MoE 的训练效率优势在推理端面临严峻的工程挑战。

\subsection{内存瓶颈与专家卸载}

MoE 模型的推理面临一个核心矛盾：
\textbf{每次只激活 5\%--10\% 的参数，
但 100\% 的参数都必须驻留在可访问的存储中}。
DeepSeek-V3 的 671B 参数在 BF16 下需要 $\sim$1.3\,TB 内存，
至少需要 17 块 H100（80\,GB）才能装下，
但每个 token 实际只需要 37B 参数（$\sim$74\,GB）的计算。

\textbf{专家卸载}（Expert Offloading）将不活跃的专家
从 GPU 显存卸载到 CPU 内存甚至 SSD，
只在需要时加载到 GPU：

\textbf{FineMoE}~(2026) 提出了\textbf{细粒度}的卸载策略：% NEW REF: finemoe2026 = FineMoE: Taming Latency-Memory Trade-Off in MoE-Based LLM Serving via Fine-Grained Expert Offloading, Yu et al., EuroSys 2026
不是以整个专家为单位进行卸载/加载
（粒度太粗——一个专家可能有几百 MB），
而是将每个专家进一步拆分为更小的“子专家”块，
按需加载。
关键洞察是：专家的使用频率遵循\textbf{长尾分布}，
少数“热门专家”被频繁调用，大部分专家很少使用。
FineMoE 将热门专家常驻 GPU，
冷门专家按子块粒度从 CPU 按需传输。

\textbf{ADEPT}（2026）进一步引入了% NEW REF: adept2026 = Two-Stage Expert Offloading for Domain-Aware MoE Inference, Kim et al., IEEE Access 2026
\textbf{领域感知的预取}策略：
根据输入文本的领域特征（代码、数学、自然语言等）
预判哪些专家可能被激活，提前将它们从 CPU 加载到 GPU。
这将专家加载操作从推理的关键路径上移除，
有效隐藏了 PCIe 总线的 I/O 延迟。

ADEPT 的预取准确率取决于一个关键假设：
\textbf{连续 token 的路由决策具有局部性}，
即当前 token 激活的专家集合
与下一个 token 的激活集合有大量重叠。
这个假设在自然语言文本中通常成立
（一段讨论物理学的文本中，
连续的 token 大概率都会激活“科学”相关的专家），
但在文本类型快速变化的场景中可能失效
（例如 interleaved code-text 格式）。
ADEPT 通过一个轻量级的 bigram 预测模型
（基于当前和前一个 token 的路由历史）
来提高预取的命中率。
实测在英文自然语言文本上
预取命中率约 85\%--92\%，
在代码混合文本上约 70\%--80\%。

专家卸载的一个实际限制是 \textbf{PCIe 带宽}。
GPU 与 CPU 之间通过 PCIe 连接
（PCIe 4.0 约 32 GB/s，PCIe 5.0 约 64 GB/s）。
加载一个 500MB 的专家从 CPU 到 GPU
在 PCIe 4.0 上需要约 15ms，
这在 prefill 阶段是可以接受的
（因为 prefill 本身就需要数百毫秒），
但在 decode 阶段（每个 token 需要约 10--50ms），
15ms 的额外延迟意味着吞吐量下降 30\%--150\%。
DeepSeek 使用的细粒度专家（每个约 0.5B 参数，$\sim$1\,GB BF16）
在这方面有优势，加载 FineMoE 的“子专家”块
可能只需 2--5ms，更容易被计算延迟所隐藏。

\subsection{MoE 模型的量化}

量化是降低 MoE 内存占用的另一条路径，
但 MoE 的量化比稠密模型更具挑战性。

\textbf{DynaExq}~(2025) 提出了% NEW REF: dynaexq2025 = Dynamic Expert Quantization for Scalable MoE Inference, Chu et al., arXiv 2025
\textbf{动态混合精度}的量化方案：
不同专家使用不同的量化位宽。
热门专家（频繁使用）保持较高精度（如 8 bit），
冷门专家使用更激进的压缩（如 4 bit 甚至 2 bit），
因为它们对最终输出的贡献权重较小。
DynaExq 在运行时根据路由统计量
动态调整专家的量化位宽，
在精度和内存之间找到最优平衡。

\textbf{MiLo}（Mixture of Low-Rank Compensators）% NEW REF: milo2025 = MiLo: Efficient Quantized MoE Inference with Mixture of Low-Rank Compensators, Huang et al., arXiv 2025
则在量化后为每个专家附加一个\textbf{低秩补偿器}，
用少量的额外参数修正量化带来的误差。
对于 4-bit 量化的 MoE 模型，
MiLo 可以恢复 90\%+ 的原始精度，
而额外的参数开销不到原始专家的 5\%。

为什么 MoE 比稠密模型更难量化？
根本原因在于\textbf{激活值的异质性}。
在稠密模型中，每一层处理所有 token，
各层的激活值分布相对稳定，
量化参数（scale 和 zero point）可以基于
所有 token 的统计特性来确定。
但在 MoE 中，不同专家只处理它们被路由到的 token 子集，
这些子集的统计特性可能差异巨大。
一个专门处理数学 token 的专家
和一个处理自然语言的专家，
它们的激活值分布、范围和形状都可能截然不同。
为所有专家使用统一的量化参数会导致某些专家的量化误差不可接受。

解决方案是\textbf{逐专家量化}（per-expert quantization）：
每个专家独立确定自己的量化参数。
DynaExq 更进一步，允许同一个专家在不同层
使用不同的量化精度，
因为同一个专家在模型底层和顶层的激活值分布也可能不同。
这种精细化的量化策略在实现上增加了复杂度，
但对最终精度的保持至关重要。

\subsection{专家并行与通信模式}

MoE 推理在多 GPU 部署时
引入了一种独特的并行模式，\textbf{专家并行}（Expert Parallelism, EP）。
与张量并行（TP，切分矩阵）和流水线并行（PP，切分层）不同，
EP 将不同的专家分布到不同的 GPU 上。

EP 的通信模式是 \textbf{All-to-All}：
每块 GPU 需要将不属于本地专家的 token 发送到远程 GPU，
同时接收其他 GPU 发来的、属于本地专家的 token。
这种通信模式与 AllReduce（所有 GPU 发送和接收相同量的数据）
有根本性差异：All-to-All 的通信量取决于路由决策，
是\textbf{动态的、不可预测的}。

在实际部署中，EP 通常与 TP 或 PP 组合使用。
一种常见的配置是在节点内使用 TP（利用 NVLink 高带宽），
在节点间使用 EP（通过 InfiniBand 传输 token）。
这种组合的挑战在于\textbf{两种并行策略的通信模式不兼容}，
TP 的 AllReduce 要求所有 GPU 同步参与，
EP 的 All-to-All 可能导致不同 GPU 之间的等待时间不同。
如何调度这两种通信以最小化总等待时间，
是 MoE 推理系统工程中最复杂的问题之一。

DeepSeek 的解决方案是将 EP 与 TP 完全分离：
在每个节点内完成 TP 的计算和通信后，
再执行跨节点的 EP All-to-All，
两者在时间上串行，避免了通信冲突。
NVIDIA 的 Megatron Core MoE~\cite{yan2026megatronmoe}
则采用了更激进的方案：
将 All-to-All 通信分解为多个小操作，
与 TP 的 AllReduce 在时间维度上交错执行，
最大化网络带宽利用率。

\subsection{批处理的矛盾}

MoE 推理还面临一个“批处理悖论”：
\textbf{批处理是推理吞吐量的关键，
但它抵消了 MoE 的稀疏性优势}。

在单请求推理中，每个 token 只激活 $K$ 个专家，
MoE 的计算节省是真实的。
但在批处理（batched inference）中，
不同请求的 token 可能路由到不同的专家，
当 batch size 足够大时，
所有 $N$ 个专家都会被至少一个 token 激活，
等价于运行一个\textbf{完整的稠密模型}。
更糟的是，MoE 的 All-to-All 通信开销在稠密模型中不存在。

\textbf{Lynx}~(2025) 提出了一种“专家重映射”方案：% NEW REF: lynx2025 = Lynx: Enabling Efficient MoE Inference via Dynamic Batch-Aware Expert Selection, arXiv 2025
在运行时动态地将 batch 中不同请求的 token
重新分配到重叠度更高的专家子集上，
减少实际需要激活的专家总数。
这在保持 $<$1\% 精度损失的前提下，
将推理吞吐量提升了 $1.23\times$。

\textbf{MoE-SpeQ}~(2025) 则将推测解码% NEW REF: moespeq2025 = MoE-SpeQ: Speculative Quantized Decoding with Proactive Expert Prefetching and Offloading for MoE, Wang et al., arXiv 2025
与专家预取结合：
在主模型验证候选 token 的同时，
利用路由器的预判结果预加载下一步可能需要的专家。
推测解码本身减少了生成步数，
专家预取则隐藏了每步的加载延迟——
两种优化相互增强。

\subsection{MoE 的优势与挑战}

\begin{corebox}[MoE vs Dense：何时选择哪种架构]
\textbf{MoE 的优势}：
\begin{itemize}[noitemsep]
  \item 知识容量与推理成本解耦——大容量、低推理成本
  \item 对于广泛部署的模型，推理经济性是压倒性优势
  \item 专家专业化带来的任务适配能力
\end{itemize}

\textbf{MoE 的挑战}：
\begin{itemize}[noitemsep]
  \item \textbf{内存占用}：所有专家必须加载到内存中，
        即使每次只激活一小部分。
        671B MoE 模型的权重（BF16）约 1.3\,TB——
        需要至少 17 块 H100（80\,GB）才能放下。
  \item \textbf{微调困难}：
        LoRA 等参数高效微调方法
        在 MoE 上的应用并不直接——
        应该给每个专家加 LoRA 适配器，
        还是只给共享部分加？
        给所有 256 个专家都加 LoRA 会导致参数量爆炸，
        只给部分专家加则可能破坏路由的平衡。
  \item \textbf{训练复杂度}：
        负载均衡、通信优化、容量因子调节——
        MoE 的超参数空间比稠密模型大得多。
  \item \textbf{推理延迟}：
        虽然 FLOP 更少，但 All-to-All 通信引入了额外延迟。
        在单 GPU 推理中，MoE 的延迟优势可能被
        更多的内存访问抵消
        （所有专家的参数都在内存中，cache 命中率更低）。
\end{itemize}

\textbf{经验法则}：
当目标模型的知识容量需要 $>$100B 参数，
且推理成本是重要考量时，选择 MoE。
当模型规模在 7B--14B，
且部署环境受限（单 GPU、边缘设备）时，选择 Dense。
\end{corebox}

\begin{whybox}[为什么 MoE 是 2026 年的标准选择？]
答案是\textbf{推理经济学}。训练一次，推理无数次。
MoE 使得你可以用 37B 的推理成本获得 671B 的知识容量。
当你服务数百万用户、每天处理数十亿 token 时，
这个 $\sim$20$\times$ 的推理计算节省是巨大的。
这就是为什么 2026 年所有 frontier 开源模型
（DeepSeek-V3/V4、LLaMA 4、Qwen3-235B、Mistral Large 3）
都采用了 MoE 架构。

值得注意的是，MoE 并不总是最优选择。
对于参数量在 7B--14B 的小模型，
Dense（稠密）架构仍然占主导——
因为小模型的参数本来就不多，
再拆分成多个专家后每个专家太小，
专业化程度不够。
而且小模型通常部署在单块 GPU 甚至 CPU 上，
MoE 的 All-to-All 通信开销在单设备上没有意义。

MoE 的真正优势在大规模模型上：
当你需要 $>$100B 的知识容量，
但推理成本必须控制在 $\sim$20B 参数以内时，
MoE 是唯一的选择。

用一个经济学类比来理解：
假设你经营一个 AI API 服务，
每天处理 10 亿个 token，
GPU 算力成本约为 \$0.001 per 1M FLOP。
一个 70B 的稠密模型每个 token 需要 $\sim$140 GFLOP，
日成本约 \$140,000。
一个 671B 的 MoE 模型（活跃 37B）
每个 token 只需要 $\sim$74 GFLOP，
日成本约 \$74,000——
虽然模型“大”了近 10 倍，推理成本反而\textbf{减半}。
年化节省约 \$24M——
这足以支付整个模型的训练成本（\$5.6M）四次。
这就是 MoE 推理经济学的核心逻辑：
\textbf{训练成本是一次性的，推理成本是持续的}，
MoE 的投资回报在持续推理中不断累积。
\end{whybox}

\section{2025--2026 年 MoE 架构的爆发}
\label{sec:moe-explosion}

2025 年下半年至 2026 年初，
MoE 从少数模型的架构选择演变为前沿中国模型的\textbf{标准范式}。
短短几个月内，至少 9 个主要模型采用了 MoE 架构，
总参数规模从 80B 跨越到 1T，
而活跃参数却压缩到惊人的低水平。

\begin{center}\small
\renewcommand{\arraystretch}{1.3}
\begin{tabular}{llrrll}
  \toprule
  \rowcolor{figLightGray}
  \textbf{模型} & \textbf{发布} & \textbf{总参数} & \textbf{活跃}
  & \textbf{专家配置} & \textbf{亮点} \\
  \midrule
  DeepSeek V4$^\star$ & 2026-04（预期） & $\sim$1T & --
  & -- & Engram + mHC \\
  Kimi K2.5 & 2026-01 & 1T & 32B
  & 384（8 act.） & Agent Swarm \\
  GLM-5 & 2026-02 & 744B & 40B
  & -- & 零 NVIDIA \\
  Qwen 3.5 & 2026-02 & 397B & 17B
  & -- & Gated DeltaNet \\
  Qwen3-Next & 2025-09 & 80B & 3B
  & 512（10+1 act.） & 超稀疏 \\
  Hunyuan 2.0 & 2025-12 & 406B & 32B
  & -- & 腾讯 \\
  MiniMax M2.5 & 2026-02 & 230B & 10B
  & -- & SWE 80.2\% \\
  Step 3.5 Flash & 2026-02 & 196B & 11B
  & -- & RTX 4090 \\
  豆包 Seed 2.0 & 2026-02 & -- & --
  & -- & AIME 98.3\% \\
  \bottomrule
\end{tabular}

\small $^\star$\,DeepSeek V4 截至 2026 年 3 月尚未正式发布，
参数为预期值。
\end{center}

这张表格揭示了 2025--2026 年 MoE 生态的三个核心趋势。

\subsection{超稀疏激活}

MoE 的稀疏度正在被推向新的极端。
DeepSeek-V3 的激活率约为 5.5\%（37B / 671B），
已经令人印象深刻——
但 \textbf{Qwen3-Next} 将这个数字压低到了
\textbf{3.75\%}（3B / 80B）。

Qwen3-Next 使用了 512 个路由专家，
每个 token 仅激活 10+1 个（10 个路由 + 1 个共享）。
80B 的模型只需要 3B 的推理计算量——
这意味着一个具有 80B 参数知识容量的模型，
推理成本仅相当于一个小型 3B 稠密模型。
激活比例从 DeepSeek-V3 的 1:18 演进到 Qwen3-Next 的 1:27，
超稀疏激活正在成为 MoE 设计的新前沿。

Qwen3-Next 的技术报告揭示了实现超稀疏激活的关键设计选择。
512 个专家中每个 token 只激活 10 个，
这意味着单次路由决策的组合空间达到了
$\binom{512}{10} \approx 2.6 \times 10^{18}$——
一个天文数字级的组合空间。
理论上，如此巨大的组合空间允许模型
为每种可能的输入模式找到一个独特的专家组合。
但实际上，路由器（一个简单的线性层）
是否有足够的表达能力来有效利用这个组合空间，
仍然是一个开放问题。

Qwen3-Next 还引入了一个有趣的设计：
10 个路由专家之外的第 11 个\textbf{共享专家}
不经过路由器选择，始终参与计算。
这意味着即使路由器的选择不够理想，
共享专家仍然为每个 token 提供了一个
高质量的“兜底”输出。
在 3.75\% 的极端稀疏度下，
共享专家的角色从“锦上添花”变成了“保命”——
它确保了每个 token 至少有一个
充分训练的专家来处理。

超稀疏激活引发了几个深层的技术问题。
首先是\textbf{稀疏度的理论极限}。
直觉上，当每个 token 只激活极少数专家时，
模型的表达能力应该在某个点开始退化——
因为复杂的推理可能需要多个专家的协同计算。
然而，当前的实验结果并未观察到明显的退化拐点：
Qwen3-Next 在 3.75\% 的激活率下
依然在多项基准上超越同规模的稠密模型。
这暗示 MoE 中的大部分专家
存在高度的功能冗余——
不同专家学到的表示之间有大量重叠。

其次是\textbf{负载均衡在超稀疏设置下的挑战}。
当只有 10/512 个专家被激活时，
路由器的偏差会被放大——
一个稍有偏好的路由器可能导致
少数专家承担绝大多数流量，
而其余专家几乎从未被激活（“专家坍缩”）。
DeepSeek 的无辅助损失方案在中等稀疏度下表现良好，
但在 $<$4\% 激活率的极端稀疏设置下
是否仍然稳定，是一个开放问题。

最后是\textbf{通信瓶颈的转移}。
超稀疏激活大幅减少了计算量，
但 All-to-All 通信模式并未等比例简化——
路由决策本身和 token 分发
仍然需要跨 GPU 通信。
当计算量降至极低水平时，
通信延迟可能反而成为推理的主要瓶颈。
DeepEP 等专用通信库正是为解决这一问题而设计的。

\subsection{消费级硬件部署}

\textbf{Step 3.5 Flash} 证明了 196B 参数的 MoE 模型
可以在\textbf{单块 RTX 4090}（24\,GB 显存）上运行。
11B 的活跃参数经过 4-bit 量化后仅需约 6\,GB，
其余参数通过专家卸载动态加载。

这打破了“大模型需要大集群”的思维定式——
MoE 的稀疏激活特性使得总参数量不再是部署的瓶颈，
活跃参数量才是。
当活跃参数降到 10--17B 范围时，
消费级 GPU 就足以提供可接受的推理体验。

消费级部署的核心机制是\textbf{专家卸载}（expert offloading）。
由于每个 token 只激活少数专家，
大部分专家参数在任意时刻都是“冷”的——
这些冷专家可以存储在更便宜的层级中：
从 GPU HBM 卸载到系统内存（DRAM），
甚至进一步卸载到 NVMe SSD。
当路由器决定激活某个被卸载的专家时，
系统通过 PCIe 或 NVMe 总线
将其参数加载回 GPU——
这引入了额外的延迟，
但由于 MoE 推理本身是内存带宽瓶颈而非计算瓶颈，
加载延迟可以与计算部分重叠。

实际部署数据表明，
在 RTX 4090 上运行 Step 3.5 Flash 时：
热专家（最常被激活的前 $k$ 个专家）驻留在 GPU HBM 中，
冷专家存储在系统 DRAM 中。
以 4-bit 量化为例，
11B 活跃参数约占 6\,GB HBM，
剩余 185B 参数约占 100\,GB 系统内存——
对于 128\,GB 内存的消费级工作站完全可行。
推理吞吐量约为数据中心部署的 1/3--1/5，
首 token 延迟在 1--3 秒范围——
对于交互式使用场景可以接受。

量化与 MoE 的交互值得特别注意。
4-bit 量化在稠密模型上已被广泛验证，
但在 MoE 场景下存在额外的风险：
量化误差可能扰动路由器的分数排序，
导致不同的专家被激活——
即\textbf{路由翻转}（routing flip）。
当前的工程解决方案是
对路由器权重保持更高精度（FP16 或 BF16），
仅对专家的 FFN 权重进行量化。

\subsection{总参数与活跃参数的分离}

从上表可以清晰看到，
行业正在收敛到一个模式：
\textbf{10--40B 的活跃参数 + 200B--1T 的总参数}。
这意味着知识容量（总参数）和推理成本（活跃参数）
被彻底解耦——
模型可以拥有巨大的知识储备，
但每次推理只调用极小一部分。

这种分离在概念上类似于人脑的工作方式：
大脑有约 860 亿个神经元，
但在任意时刻只有极小比例同时活跃。
MoE 的总参数对应大脑的总突触连接数（长期记忆容量），
活跃参数对应工作记忆中同时处理的信息量。
这个类比虽然粗糙，
但揭示了一个可能的设计原则：
\textbf{知识存储和在线计算有各自独立的缩放规律}。

活跃参数量的收敛值（10--40B）也很有启发性。
它暗示了在当前的 token 表示维度（$d = 4096$--$8192$）
和任务复杂度下，
每个 token 需要的“同时在线计算能力”
大约在 10--40B 参数对应的 FLOP 范围内。
低于这个范围（如 Qwen3-Next 的 3B），
虽然仍然有效，但在复杂推理任务上可能受限——
因为每个 token 能调动的专家过少，
协同计算能力不足。
高于这个范围（如 LLaMA 4 Behemoth 的 288B），
则可能出现算力浪费——
大量活跃参数中有相当比例
在处理简单 token 时是多余的。
DynaMoE 的动态激活策略正是为了解决这种浪费——
让简单 token 用 10B 级别的计算，
让复杂 token 用 40B+ 级别的计算。

从工程角度看，
总参数与活跃参数的分离
对 serving 基础设施提出了新要求。
1T 总参数意味着约 2\,TB 的模型权重（FP16），
即使以 4-bit 量化存储也需要约 500\,GB——
这远超单台服务器的 GPU 显存容量，
需要多机分布式部署或专家卸载。
模型的\textbf{冷启动时间}
（从磁盘加载全部参数到可服务状态）
可能需要数分钟——
这对弹性伸缩（autoscaling）构成了挑战。
解决方案包括预热池（warm pool）、
专家参数的预加载策略、
以及按需加载的 lazy loading 架构。

DeepSeek V4（预期）的具体活跃参数量尚未公布，
但据传其设计目标包括原生多模态理解与生成、
Engram 长期记忆和 mHC（流形约束超连接）等创新架构。
活跃参数量的选择最终取决于
模型需要\textbf{同时在线}的能力复杂度——
多模态推理和长期记忆检索
可能需要比纯文本推理更多的活跃参数。
一个开放的研究问题是：
总参数的增长是否存在收益递减的拐点？
当总参数从 1T 增长到 10T 时，
模型的知识容量是否仍然线性增长，
还是因为专家间的知识冗余
而出现边际收益递减？

\section{MoE 与稠密模型的缩放行为对比}
\label{sec:moe-scaling-comparison}

MoE 和稠密模型遵循不同的缩放规律（scaling laws），
理解这些差异对于做出正确的架构选择至关重要。

\subsection{Chinchilla 定律在 MoE 中的适用性}

Chinchilla 定律~\cite{hoffmann2022chinchilla}
指出稠密模型的最优训练策略是
模型参数 $N$ 和训练 token 数 $D$ 的比例
应约为 $D \approx 20N$。
但这个比例是基于“所有参数每个 token 都参与计算”的假设——
在 MoE 中，这个假设不成立。

对于 MoE 模型，需要区分两个参数量：
总参数 $N_{\text{total}}$ 和每个 token 的活跃参数 $N_{\text{active}}$。
直觉上，训练 token 数应该与\textbf{活跃参数}而非总参数匹配——
因为每个 token 只“看到”活跃参数的计算。
这意味着 MoE 模型的最优训练 token 数
应该远少于等量总参数的稠密模型。

DeepSeek-V3 的训练数据量（14.8T token）与活跃参数（37B）的比例
约为 400:1——远高于 Chinchilla 的 20:1。
这看起来似乎违反了 Chinchilla 定律，
但考虑到 671B 的总参数量，
每个专家的有效训练 token 数
大约是 $14.8T \times (8/256) \approx 462B$——
与 Chinchilla 为 37B 活跃参数推荐的
$37B \times 20 = 740B$ 大致在同一数量级。
这暗示 MoE 的缩放定律中，
训练充分性应该用“每个专家的有效训练 token 数”
而非“总训练 token 数”来衡量。

\subsection{MoE 的“免费午餐”与其代价}

MoE 经常被描述为提供了“免费的”参数增长——
增加更多专家不增加推理成本（活跃参数不变），
但增加了模型的知识容量（总参数增加）。
这是否意味着我们应该无限增加专家数量？

答案是否定的。增加专家数量有三个递减收益的来源。

\textbf{知识冗余}。
当专家数量足够多时，
不同专家学到的表示之间开始出现大量重叠。
256 个专家中可能有 20 个都学会了“将英文翻译为法文”，
因为路由器无法完美地将所有翻译相关的 token
集中分配到同一个专家。
这种冗余意味着新增专家带来的边际知识增量
随专家数量的增加而下降。

\textbf{路由精度瓶颈}。
路由器是一个简单的线性层，
当专家数量从 256 增加到 1024 时，
路由器需要在 1024 个选项中做出准确的选择——
这对一个没有非线性激活函数的线性投影来说
是一个越来越困难的任务。
高维空间中的“角度集中”现象
使得不同专家的路由分数趋于相似，
路由决策变得不稳定。

\textbf{通信成本}。
更多的专家意味着更多的 GPU 上分布着专家，
All-to-All 通信的参与方增多，
通信延迟和带宽需求增加。
在某个临界点之后，
通信开销的增长可能超过新增专家带来的质量提升。

经验上，当前的实践收敛到了一个“甜蜜点”：
\textbf{128--512 个专家}，
\textbf{Top-8 路由}，
\textbf{活跃参数 20--40B}，
\textbf{总参数 200B--1T}。
超过这个范围的收益递减程度
是一个活跃的研究问题——
Kimi K2 的 384 专家和 DeepSeek 可能的下一代模型
将提供更多关于专家数量上限的实证数据。

\section{MoE 理论的新进展}
\label{sec:moe-theory-2026}

伴随 MoE 在实践中的爆发，
理论层面也涌现了重要的新成果。
这些工作从三个互补的角度推进了对 MoE 的理解：
缩放定律的理论刻画、
计算效率的极限逼近、
以及系统层面的全栈优化。

\subsection{MoE 的缩放定律：专家与注意力的最优分配}

传统 Scaling Laws（如 Chinchilla 定律）
将模型视为同质的参数集合——
“增加 $N$ 倍参数，性能提升 $f(N)$”。
但 MoE 模型打破了这种同质性假设：
模型中存在两种本质不同的计算单元，
它们对性能的贡献方式截然不同。

\textbf{Optimal Expert-Attention Allocation}% NEW REF: expertattention2026 = Optimal Expert-Attention Allocation in MoE, arXiv:2603.10379, 2026
（arXiv:2603.10379，2026 年 3 月）
首次将神经 Scaling Laws 扩展到了异构的 MoE 架构。
论文将 MoE 的计算预算分解为两个可独立调节的维度：
\textbf{专家层的容量}（负责知识检索和领域特化）和
\textbf{注意力层的容量}（负责上下文交互和信息整合）。
核心研究问题是：在固定的总计算预算 $C$ 下，
分配多少给专家层、多少给注意力层，
才能使语言建模 loss 最小化？

论文通过大规模的消融实验
（在 1B 到 100B 参数规模上系统性地扫描分配比例）
得到了几个反直觉的结论。
首先，\textbf{最优分配比例与模型规模相关}——
小模型应将更大比例的计算分配给注意力层
（因为知识容量有限，上下文理解更为稀缺），
而大模型则应将更多资源分配给专家层
（因为注意力层已经足够强，
知识存储成为瓶颈）。
其次，最优比例也与\textbf{任务类型}密切相关——
知识密集型任务（如百科问答）
偏好更多的专家容量，
而推理密集型任务（如数学证明）
则偏好更强的注意力层。
这意味着一个通用模型的最优设计
必须在不同任务需求之间寻找折中。

这些发现对 MoE 架构设计有直接的指导意义：
盲目增加专家数量（如从 256 增加到 1024）
可能不如将部分计算预算转移到更强的注意力层——
尤其是当目标任务需要复杂推理而非纯知识检索时。
这也解释了为什么 DeepSeek-V3 的注意力层
使用了 MLA（Multi-head Latent Attention）
这种相对“重”的注意力设计——
它可能正是为了平衡专家层与注意力层之间的计算分配。

\subsection{LatentMoE：潜在空间中的路由与计算}

标准 MoE 的路由在完整的 $d$-维隐空间中进行：
路由器是一个 $[d, N_{\text{experts}}]$ 的线性投影，
其中 $d$ 通常在 4096--8192 之间。
当专家数量 $N$ 很大时（如 256 或 384），
路由矩阵本身的参数量就达到百万级，
而且高维空间中的路由决策可能受到“维度诅咒”——
token 之间的相似度在高维空间中变得不那么有区分力。

\textbf{LatentMoE}（arXiv:2601.18089，2026 年 1 月）% NEW REF: latentmoe2026 = LatentMoE: Toward Optimal Accuracy per FLOP and Parameter in MoE, arXiv:2601.18089, 2026
提出了一种在\textbf{低维潜在空间}中同时执行路由和专家计算的范式。
具体来说，LatentMoE 首先将输入 $\mathbf{x} \in \mathbb{R}^d$
投影到一个低维潜在空间 $\mathbf{z} \in \mathbb{R}^r$（$r \ll d$），
在这个压缩空间中做路由决策和部分专家计算，
然后将结果投影回原始维度。

这种设计的理论优势体现在两个层面。
在\textbf{参数效率}方面，
低维空间中的路由矩阵从 $O(d \cdot N)$ 压缩到 $O(r \cdot N)$，
专家的 FFN 计算也在低维空间中进行，
参数量大幅减少。
在\textbf{路由质量}方面，
低维空间中的路由决策更加稳定——
压缩操作天然起到了降噪效果，
减少了高维空间中因“角度集中”而导致的路由振荡。
论文证明了在温和的假设下，
LatentMoE 与标准 MoE 具有相同的\textbf{通用逼近能力}，
即潜在空间中的计算不会损失表达能力。

LatentMoE 与 DeepSeek-V3 的 MLA 在设计理念上有深刻的共鸣：
两者都在核心计算之前引入了一个“瓶颈”投影，
将高维输入压缩到低维空间。
MLA 压缩的是注意力的 KV 表示，
LatentMoE 压缩的是专家的输入表示。
如果将两者组合——
注意力层用 MLA 压缩，FFN/MoE 层用 LatentMoE 压缩——
整个模型的计算过程都将在低维空间中进行，
只在输入和输出端做维度转换。
这种“全程低维计算”架构的可行性和收益
是一个值得探索的方向。

从实践角度看，LatentMoE 为超稀疏 MoE 模型
（如 Qwen3-Next 的 512 专家配置）
提供了一条有吸引力的优化路径：
在保持模型知识容量不变的前提下，
通过在潜在空间中操作，
进一步压缩每个 token 的推理成本。
初步实验显示，在 $r = d/4$ 的配置下，
LatentMoE 可以在几乎不损失精度的前提下
将路由相关的计算量减少约 60\%。

\subsection{Megatron Core MoE：系统层面的全栈优化}

前两项工作关注的是“模型应该如何设计”，
而 \textbf{NVIDIA Megatron Core MoE}（arXiv:2603.07685，2026 年 3 月）% NEW REF: yan2026megatronmoe
则回答了一个同样关键的问题：
“设计好的 MoE 模型如何高效地训练？”

大规模 MoE 训练面临一个独特的“三方耦合”困境。
增大专家数量会同时增加内存占用（更多专家权重需要存储）、
通信量（All-to-All 传输更多的 token-专家映射）、
以及计算模式的碎片化（更细粒度的专家导致 GPU 利用率下降）。
这三个维度的交互使得
\textbf{单独优化任一维度可能恶化其他维度}——
例如，通过增加流水线并行来分散内存压力，
可能引入更多的通信同步点；
通过增大 batch size 来提高计算效率，
可能导致 All-to-All 通信量暴增。

Megatron Core MoE 通过全栈协同设计
将专家并行（EP）、张量并行（TP）和流水线并行（PP）
纳入统一的抽象层。
系统的核心创新包括：
（1）\textbf{自适应并行策略选择器}——
根据模型配置（专家数量、Top-K、专家大小）
和硬件拓扑（GPU 数量、NVLink/IB 带宽）
自动搜索最优的并行方案组合；
（2）\textbf{通信--计算重叠调度器}——
将 All-to-All 通信分解为多个小操作，
与专家计算在时间维度上精细地交错执行，
最大化 GPU 利用率；
（3）\textbf{内存感知的激活重算}——
根据剩余显存动态决定哪些激活值需要重算，
在内存和计算之间找到最优权衡。

在 NVIDIA DGX H100 集群上的实测显示，
Megatron Core MoE 在 256 GPU 上训练
DeepSeek-V3 规模的模型（256 专家 + Top-8）时，
相比朴素的并行策略实现了约 $1.8\times$ 的吞吐提升，
同时将峰值显存占用降低了约 30\%。
这些数字意味着同样的硬件预算
可以训练更大的 MoE 模型、
或在更短的时间内完成训练——
对于动辄数百万美元的训练成本来说，
系统优化的经济价值不亚于算法创新。

这三项理论和系统工作共同描绘了 MoE 研究的前沿图景：
缩放定律告诉我们\textbf{如何分配}计算资源，
LatentMoE 告诉我们\textbf{如何压缩}计算过程，
Megatron Core MoE 告诉我们\textbf{如何执行}大规模训练。
三者的结合为下一代万亿参数 MoE 模型的设计和训练
提供了理论指导和工程基础。

\section{MoE 与其他稀疏架构的关系}
\label{sec:moe-related-architectures}

MoE 并非实现稀疏计算的唯一途径。
理解 MoE 在更广泛的稀疏架构谱系中的位置，
有助于把握这一领域的发展脉络。

\subsection{MoE vs 稀疏注意力}

MoE 和稀疏注意力分别在 Transformer 的两个核心组件上
引入了稀疏性。
MoE 稀疏化了\textbf{FFN 层}——
不是所有参数都参与计算，而是动态选择子集。
稀疏注意力（如 NSA、Double-P）稀疏化了\textbf{注意力层}——
不是所有 token 对都计算注意力分数，
而是只关注最相关的 token 子集。

两者的优化目标不同但互补。
MoE 减少了\textbf{每个 token 的参数计算量}，
稀疏注意力减少了\textbf{token 之间的交互计算量}。
在超长上下文（如 1M token）场景下，
注意力的 $O(n^2)$ 开销远超 FFN 的 $O(n \cdot d_{\text{ff}})$ 开销——
此时稀疏注意力的节省更为关键。
但在短上下文（如 4K token）场景下，
FFN 的计算量占总计算的大部分（约 2/3），
MoE 的节省更为显著。

这解释了为什么最先进的模型同时使用两种稀疏性：
DeepSeek-V3 在 FFN 层使用 MoE，
同时在注意力层使用 MLA 来压缩 KV cache。
LLaMA 4 在 FFN 层使用 MoE，
在注意力层使用 iRoPE（交替 NoPE 层实质上是一种
对部分层跳过位置编码计算的“稀疏”策略）。

\subsection{MoE vs 早期退出}

另一种稀疏计算策略是\textbf{早期退出}（early exit）：
对于“简单”的 token，在经过少数几层后就输出结果，
跳过后面的层。
这与 MoE 的“水平稀疏”（选择专家子集）形成互补——
早期退出是“垂直稀疏”（选择层子集）。

两者的结合——“同时选择哪些层和哪些专家来处理每个 token”——
是一个尚未被充分探索的设计空间。
DynaMoE~\cite{gulme2026dynamoe} 的层间可变容量设计
在这个方向上迈出了一步：
通过为不同层分配不同的专家容量，
它隐式地实现了一种“部分早期退出”的效果——
对于不需要深层处理的 token，
后续层激活更少的专家（等效于降低这些层的计算强度），
虽然 token 仍然通过了所有层，
但实际计算量因为专家激活减少而降低。

\subsection{MoE vs 状态空间模型}

SSM（State Space Model，如 Mamba）
代表了另一种效率范式：
通过线性递推替代注意力机制，
将序列处理从 $O(n^2)$ 降低到 $O(n)$。
MoE 和 SSM 可以组合使用——
Jamba 模型~\cite{jamba2024}（AI21 Labs）% NEW REF: jamba2024 = Jamba: A Hybrid Transformer-Mamba Language Model, AI21 Labs, arXiv 2024
就是第一个将 Transformer 注意力层、Mamba 层和 MoE 层
混合在同一个模型中的架构。

在 Jamba 的设计中，
部分层使用 Mamba 的线性递推（高效处理长序列），
部分层使用 Transformer 注意力（精确的长距离检索），
两者的 FFN 部分都可以替换为 MoE。
这种“三元混合”架构的设计空间极其庞大——
哪些层用 Mamba、哪些用注意力、哪些用 MoE
以及它们的组合模式，
需要大量的实验来确定最优配置。
Jamba 的初步结果表明，
这种混合架构在保持注意力精度的同时，
显著降低了长序列的处理成本——
一个 52B 参数的 Jamba 模型
在 256K 上下文上的推理吞吐量
约为同规模纯 Transformer + MoE 模型的 $1.5\times$。

\subsection{MoE 的未来：走向自适应计算}

回顾 MoE 的发展轨迹，
一个清晰的主线是\textbf{计算的自适应性}在不断增强。

从 Switch Transformer 的固定 Top-1
到 DynaMoE 的动态 Top-K，
从所有层统一的专家配置
到层间可变的容量分配，
从固定精度的推理
到 DynaExq 的动态混合精度——
每一步都在让模型更智能地分配自己的计算资源。

终极目标可能是一种完全自适应的计算架构：
\textbf{每个 token 在每一层
自动决定需要多少计算、
使用哪些专家、
以什么精度执行，
甚至是否需要通过这一层}。
这不仅是 MoE 的进化方向，
也是更广泛的“高效计算”研究的核心愿景。
当计算资源（GPU 算力、内存、带宽）
成为 AI 系统的主要成本时，
让每一个 FLOP 都花在刀刃上——
而不是浪费在“简单 token 不需要的专家计算”
或“不需要看到全部上下文的注意力计算”上——
就是降低成本、提升可及性的关键路径。

MoE 在 2026 年的地位可以用一句话概括：
\textbf{它是大模型通向实际部署的必经之路}。
当推理成本决定了 AI 产品的可行性时，
能够用 37B 的推理成本提供 671B 的知识容量——
这不是一个学术上的优雅，而是一个商业上的必需。

%% ============================================================
% NEW REF: jacobs1991moe = Adaptive Mixtures of Local Experts, Jacobs et al., Neural Computation 1991
% NEW REF: lepikhin2021gshard = GShard: Scaling Giant Models with Conditional Computation, Lepikhin et al., ICLR 2021
% NEW REF: zhou2022expertchoice = Mixture-of-Experts with Expert Choice Routing, Zhou et al., NeurIPS 2022
%% --- NEW REFERENCES (March 2026 update) ---
% NEW REF: gulme2026dynamoe = DynaMoE: Dynamic Token-Level Expert Activation with Layer-Wise Adaptive Capacity, Guelmez, arXiv 2026
% NEW REF: dirmoe2026 = DirMoE: Differentiable Probabilistic Router for Mixture-of-Experts, ICLR 2026
% NEW REF: l2r2026 = L2R: Low-Rank and Lipschitz-Controlled Routing for Mixture-of-Experts, Yang et al., arXiv 2026
% NEW REF: llama4 = The Llama 4 Herd, Meta AI, 2025
% NEW REF: qwen3 = Qwen3 Technical Report, Qwen Team, arXiv 2025
% NEW REF: kimiTeam2025kimik2 = Kimi K2: Open Agentic Intelligence, Kimi Team, arXiv 2025
% NEW REF: finemoe2026 = FineMoE: Taming Latency-Memory Trade-Off in MoE-Based LLM Serving, Yu et al., EuroSys 2026
% NEW REF: adept2026 = Two-Stage Expert Offloading for Domain-Aware MoE Inference, Kim et al., IEEE Access 2026
% NEW REF: dynaexq2025 = Dynamic Expert Quantization for Scalable MoE Inference, Chu et al., arXiv 2025
% NEW REF: milo2025 = MiLo: Efficient Quantized MoE Inference with Low-Rank Compensators, Huang et al., arXiv 2025
% NEW REF: lynx2025 = Lynx: Enabling Efficient MoE Inference via Dynamic Batch-Aware Expert Selection, arXiv 2025
% NEW REF: moespeq2025 = MoE-SpeQ: Speculative Quantized Decoding with Expert Prefetching for MoE, Wang et al., arXiv 2025
% NEW REF: expertattention2026 = Optimal Expert-Attention Allocation in MoE, arXiv:2603.10379, 2026
% NEW REF: latentmoe2026 = LatentMoE: Toward Optimal Accuracy per FLOP and Parameter in MoE, arXiv:2601.18089, 2026
% NEW REF: yan2026megatronmoe = NVIDIA Megatron Core MoE: Full-Stack System Co-Design for MoE Training, NVIDIA, arXiv:2603.07685, 2026
% NEW REF: puigcerver2024softmoe = From Sparse to Soft Mixtures of Experts, Puigcerver et al., arXiv:2308.00951, 2023
% NEW REF: dbrx2024 = DBRX: A New State-of-the-Art Open LLM, Databricks/Mosaic, 2024
% NEW REF: hoffmann2022chinchilla = Training Compute-Optimal Large Language Models, Hoffmann et al., NeurIPS 2022
% NEW REF: jamba2024 = Jamba: A Hybrid Transformer-Mamba Language Model, AI21 Labs, arXiv 2024
